%% Plantilla LaTeX Espacios Editorial UD
\documentclass[draft,12pt,twoside,spanish]{book}

\usepackage{verbatim}
\usepackage{listings} % Para codigo fuente y trozos de codigo
\lstloadlanguages{C,bash}
\lstset{language=C, 
        basicstyle=\scriptsize\ttfamily,
        breaklines=true,
        keywordstyle=\color{blue}\ttfamily,
        stringstyle=\color{red}\ttfamily,
        commentstyle=\color{gray}\ttfamily,
        morecomment=[l][\color{magenta}]{\#},
        numbers=left, numberstyle=\tiny
}
\usepackage{minted}
\usepackage[spanish,es-tabla]{babel}
\usepackage{siunitx}
\usepackage{url}
\usepackage{graphicx}
\graphicspath{{capitulos/figuras/}}
\usepackage[]{caption} %labelsep=period
\usepackage{float}

\usepackage[
backend=biber,
style=apa,
]{biblatex}
%%% Editar %%%
\addbibresource{references.bib}
\let\cite\parencite
%%%%%%%%%%%%%%

\usepackage{fontspec}
\input{esConfigDir/esConfigFonts}
\usepackage{authblk}
\DefineBibliographyStrings{spanish}{andothers={\itshape{et~al}\adddot}}

%In a current version of biblatex-apa, the comma before the "and" is either by the Oxford comma (\finalnamdelim)
%or produced by the internal macro \apablx@ifrevnameappcomma.
%If you want to get rid of the comma, you need to disable both of these features.
%https://tex.stackexchange.com/questions/630515/biblatex-remove-commas-after-author-first-names-in-references#:~:text=1%20Answer&text=In%20a%20current%20version%20of,disable%20both%20of%20these%20features.

\makeatletter
\renewcommand*{\apablx@ifrevnameappcomma}[2]{#2}
\makeatother
\DefineBibliographyExtras{spanish}{%
  \let\finalandcomma=\empty
}
%Para cambiar el & por "y" en el capitulo de Referencias
\DeclareDelimFormat[bib,textcite]{finalnamedelim}{%
  \ifnumgreater{\value{liststop}}{2}{}{}%
  \addspace\bibstring{and}\space}

%para cambiar & en las citaciones dentro del texto
\DeclareDelimFormat[parencite]{finalnamedelim}{\addspace{y}\space}
\setcounter{biburlbreakpenalty}{20}

\usepackage[letterspace=5]{microtype}
\usepackage{graphicx}

%%% Editar %%%
\def\texttitulocorto{Evaluación de una red IoT bajo el estándar LoRa}
%Lista de autores
\def\loauthors{J.\,N.~Poveda, L.\,M.~Santamaría, E.\,E.~Gaona García, S.\,Tabares, A.\,F.~Bustos}
%%%%%%%%%%%%%%
\input{esConfigDir/esEstilo}

\usepackage[
  paperheight=24.0cm,
  paperwidth=17.0cm,
  textwidth=12.5cm,
  tmargin=2.5cm,
  bmargin=2.5cm,
  left=2.5cm,
  right=2.0cm
]{geometry}


%%% Editar %%%
%Algunos paquetes para el usuario
\decimalpoint
\usepackage{lipsum}  
\usepackage[pangram]{blindtext}
%%%%%%%%%%%%%%


%%% Editar %%%
%% Comentar esta línea para la presentación si hipertexto
\usepackage[unicode,bookmarksopen,bookmarksdepth=2,breaklinks=true]{hyperref}%Para hipervinulos de figuras, tablas, ecuaciones y citas
\hypersetup{
        colorlinks=true,
        linkcolor=black,
        filecolor=black,
        urlcolor=black,
        citecolor=black
    }
\usepackage{etoolbox}
\gappto{\UrlBreaks}{\UrlOrds}
%%%%%%%%%%%%%%

\usepackage{amsmath}
\usepackage{longtable}
\usepackage{subcaption}
\usepackage{array}
\newcolumntype{P}[1]{>{\centering\arraybackslash}p{#1}}

\usepackage{shapepar}
\usepackage{tikz}
\usepackage{tikzscale}
\usepackage{xcolor}
\definecolor{orcidlogocol}{HTML}{A6CE39}

\usepackage{tabularray}
\usepackage{nameref}
\usepackage{fontawesome}
\usepackage{academicons}

%% Para editar: comentar esta línea.
%% Para produccion del libro: activar esta linea.
%% Nota: Esta instrucción es para colocar lineas guia de corte (crop marks)
%\usepackage[cam,width=190truemm,height=260truemm,lualatex,center]{crop}
%%%%%%%%%%

%%% Editar %%%
%%%%%%%%%%%%%%

\usepackage[acronym,toc]{glossaries}
\usepackage{amsmath,amssymb,amsfonts,latexsym,cancel}
\usepackage{multirow}
\usepackage{booktabs}
\usepackage[table]{xcolor}
\usepackage{tikz} % hacer gráficos 
\usepackage{setspace}

\def\fuente#1{\small\textbf{Fuente:}~#1}
\date{\today}

\newglossarystyle{long-tabla}{%
  \setglossarystyle{long}

  \renewcommand{\glossentry}[2]{%
    \small
    \glsentryitem{##1}\glstarget{##1}{\textmed{\glossentryname{##1}}} &
    \textls[-10]{\Glossentrydesc{##1}}\glspostdescription\space ##2\tabularnewline
  }%
}

\newglossarystyle{long-lista}{%
  \setglossarystyle{list}
  \renewcommand*{\glsgroupskip}{}%
  \renewcommand*{\glossentry}[2]{%
    \small
    \item[\glstarget{##1}{\textmed{\glossentryname{##1}}}] 
    \textls[-10]{\Glossentrydesc{##1}}\glspostdescription\space ##2
  }%
}

\setglossarystyle{long-lista}

%%% Sin sangria al inicio de párrafos
\setlength\parindent{0pt}
\raggedbottom
%%%

\pagestyle{empty}
\counterwithout{footnote}{chapter}

\makeglossaries

\def\texttitulo{Evaluación de una red IoT bajo el estándar LoRa}
\def\textsubtitulo{}
\def\loauthors{José~Noé~Poveda~Zafra, Luis~Enrique~Martín~Santamaría, Elvis~Eduardo~Gaona~García, Sebastián~Yesid~Tabares~Amaya, Alex~Francisco~Bustos~Pinzón}
\def\texttitulocorto{\texttitulo}

\usepackage{emptypage}

% Esto asegura que la página izquierda de un capítulo nuevo esté vacía (sin cornisas)
\renewcommand{\cleardoublepage}{%
  \clearpage\ifodd\value{page}\else\hbox{}\thispagestyle{empty}\newpage\fi%
}

\usepackage[all]{nowidow}

\newcommand{\CO}{CO\textsubscript{2}}
\newcommand{\IC}{I\textsuperscript{2}C}

%%%%%%%%%%%%%%%%%%%%%%%%%%%%%%%%%%%
%%%
%%% Begin document
%%%
%%%%%%%%%%%%%%%%%%%%%%%%%%%%%%%%%%%

\begin{document}

\renewcommand*\contentsname{\lhseries Contenido}
\renewcommand{\cfttoctitlefont}{\hfill\Huge}
\renewcommand{\cftdot}{}
\renewcommand*{\bibname}{\fontschap Referencias}

\cleardoublepage
\input{portada.tex}
\cleardoublepage
\input{portada_dos}
\cleardoublepage

\abstract
{Resumen}
{El libro analiza la implementación y evaluación de una red IoT basada en el estándar LoRaWAN, destacando su eficiencia para la transmisión de datos en aplicaciones de monitoreo ambiental. Además, describe la arquitectura de la red, incluyendo nodos, pasarelas y servidores, así como sus capas (física, MAC y aplicación). Igualmente, se detalla la selección e integración de sensores para medir la calidad del aire y el uso de \emph{software} libre para la gestión y visualización de datos. A través de un enfoque experimental, se diseñó, simuló e implementó la red en Bogotá, evaluando cobertura, interoperabilidad y rendimiento. Los resultados evidencian la viabilidad de LoRaWAN para soluciones escalables, de bajo consumo y aplicables a ciudades inteligentes.}
{Palabras clave}
{IoT, LoRaWAN, sensores, calidad del aire, redes inalámbricas, monitoreo.}

\noindent\rule{1.8cm}{0.4pt}

\vspace{3mm}

\begin{otherlanguage*}{english}
  \abstract
  {Abstract}
  {The book analyzes the implementation and evaluation of an IoT network based on the LoRaWAN standard, highlighting its efficiency for data transmission in environmental monitoring applications. It also describes the network architecture, including nodes, gateways, and servers, as well as its layers (physical, MAC, and application). Likewise, it details the selection and integration of sensors to measure air quality and the use of open-source \emph{software} for data management and visualization. Through an experimental approach, the network was designed, simulated, and implemented in Bogotá, evaluating coverage, interoperability, and performance. The results demonstrate the viability of LoRaWAN for scalable, low-power solutions applicable to smart cities.}
  {Keywords}
  {IoT, LoRaWAN, sensors, air quality, wireless networks, monitoring.}
\end{otherlanguage*}

\vspace{1cm}

\hfill\comocitar{Poveda Zafra,~J.\,N., Martín Santamaría,~L.\,E., Gaona García,~E.\,E., Tabares Amaya,~S.\,Y. y Bustos Pinzón,~A.\,F. (2026). Evaluación de una red IoT bajo el estándar LoRa. Universidad Distrital Francisco José de Caldas-Editorial UD.}{XXXXXXXXXX}


\cleardoublepage

\chapteralt{Agradecimientos}

Los autores agradecen a la Oficina de Investigaciones de la Universidad Distrital Francisco José de Caldas por el apoyo en la financiación del proyecto de investigación titulado ``Evaluación de una red IoT de rango amplio bajo el estándar LoRa'', que permitió la publicación de los resultados obtenidos en el proyecto desarrollado en el grupo de investigación GITUD.


%% Editar re-definiciones de comandos para el usuario%%
%% panga aquí sus definiciones
%%
%%
%%%%%%%%%%%%%%
\setcounter{tocdepth}{1}
\tableofcontents
\addtocontents{toc}{\protect\thispagestyle{empty}}
\cleardoublepage

\pagestyle{esstyle}

\chapter{Introducción}

El propósito de este libro de investigación es analizar en profundidad los modelos de \gls{iot} abiertos y comprender cómo pueden proporcionar soluciones efectivas para la lectura y transmisión de datos que los objetos ofrecen a través de los sensores. Al comprender estos modelos, será posible controlar o leer parámetros de los objetos para aplicaciones inteligentes. Es importante destacar que el uso de sistemas de hardware y software abiertos ofrece numerosas ventajas, como la interoperabilidad entre diferentes fabricantes y el acceso a las arquitecturas y códigos necesarios para ajustar los proyectos de \gls{iot} a las necesidades específicas. Por lo tanto, este libro tiene como objetivo brindar a los lectores una comprensión clara de los modelos de IoT abiertos y su aplicación en diferentes campos.

Este libro es el resultado de un proyecto de investigación en el que se ha desplegado una red \gls{lora}. El objetivo del proyecto es desarrollar un sistema de nodos, concentradores y servidores, junto con software libre diseñado para la gestión de la red y el tratamiento de datos. El objetivo final es proporcionar resultados eficaces, desde la recolección de los datos hasta la entrega de la información a través de una interfaz web.

El libro está estructurado en capítulos que abordan temas específicos de la red \gls{lora}. En los capítulos 2 y 3,  se describen las generalidades y los los fundamentos de la red \gls{lora}, desde los protocolos hasta la descripción de cada capa de la red. Se comienza con la capa física, seguida de la capa de MAC, que se despliega en tres clases: Clase A, Clase B y Clase C. Por último, se describe la capa de aplicación. Se proporcionará una descripción detallada de las tramas de recepción y transmisión que involucran cada una de las capas de la red \gls{lora}. También se incluirá una descripción de las bandas de frecuencia asignadas a este estándar, así como los tipos de modulación y demodulación asignados a la transmisión y recepción. La característica más relevante de la red \gls{lora} es la seguridad de los datos, que es vital en el despliegue radioeléctrico de la información.

En el capítulo 4, se abordarán las arquitecturas de los nodos propuestos por el estándar \gls{lora}. Se describirán detalladamente cada uno de los bloques que componen el nodo. En primer lugar, se analizará la fuente de alimentación y se especificarán los márgenes de tolerancia requeridos para su correcto funcionamiento. A continuación, se estudiará la inteligencia del nodo a través de un microcontrolador que deberá cumplir con ciertos criterios para ser utilizado en el proyecto, como la capacidad de adquirir datos, el reloj, la gestión de los protocolos del nodo, los periféricos para la conexión de los sensores y la capacidad de procesamiento de datos, entre otros.

También se abordará el bloque de radio \gls{lora}, que se caracteriza por su bajo consumo de energía y se mostrará su arquitectura y el manejo de los periféricos. En este capítulo se destacará el diseño del nodo para el proyecto, que puede implementarse basado en un chipset \gls{lora}, un RF-MCU o un módem \gls{lora} interno o externo. Asimismo, se describirá la interfaz digital del nodo, la cual debe ser compatible con el estándar \gls{lora} y utilizar las principales interfaces como $I^2C$, SPI y UART. Se presentará una descripción detallada de la interfaz utilizada en el proyecto, así como la interfaz de hardware con la disposición sugerida para su correcto funcionamiento.

En el capítulo 5, además de introducir los sensores, se profundizará en la selección de los mismos para el proyecto en cuestión. Se explicará detalladamente cada uno de los sensores necesarios y su respectiva interfaz, desde los sensores de interfaz analógica, que miden señales analógicas como la temperatura o la humedad, hasta los sensores de interfaz digital que permiten la comunicación directa con el microcontrolador a través de protocolos como $I^2C$, SPI, entre otros.

En el caso específico del sensor de CO2, se describirán sus características y su interfaz de despliegue, así como también se hablará de la forma en que se medirá su concentración en el ambiente y cómo se utilizarán los datos recolectados para las aplicaciones inteligentes. 

Además, se abordará el tema de los sensores de geolocalización, explicando su funcionamiento y su importancia en el proyecto, así como también se describirán otros sensores utilizados en el proyecto, como pueden ser los sensores de luminosidad, los sensores de temperatura, los sensores de humedad, entre otros. En general, se explicará el uso de cada uno de los sensores y se describirán detalladamente sus características técnicas y su interfaz de conexión.

El capítulo 6 se enfocará en el concentrador LoRa, donde se describirán detalladamente los componentes que lo conforman, como el host, el núcleo y las funcionalidades básicas para su correcto funcionamiento. Como host se utilizará el minicomponente Raspberry Pi 3, que cuenta con una interfaz GPIO para la conexión de periféricos externos. La pasarela utilizada es la encargada de la interfaz de radio y se encarga de comunicarse con miles de sensores y a través del host, llevar esos datos a un servidor LoRaWAN para la gestión de los datos.

Se empleará el sistema operativo Alpine Linux, debido a que es muy liviano y está orientado a la seguridad, además de ser de propósito general. Se presentará una guía ilustrada de cómo configurar el sistema operativo para optimizar su desempeño. Se utilizará el software LoRa Packet Forwarder para enviar los paquetes de datos al servidor, y el software LoRa Server será el encargado de gestionar las tramas de datos, tanto descendentes como ascendentes, y en general, de la programación de la transmisión de los datos de enlace.

En cuanto al almacenamiento y respaldo de datos, se utilizará el software InfluxDB como base de datos. Por último, para la visualización de los datos, se empleará el software Grafana, el cual permite a los usuarios visualizar los datos en línea y en tiempo real. Todo esto se explicará detalladamente en el capítulo 4, incluyendo el proceso de instalación y configuración de cada uno de los componentes, así como la interacción entre ellos para asegurar el correcto funcionamiento del concentrador LoRa.

En el capítulo 7, se presenta la implementación de la red \gls{lwan} en su totalidad, una vez que se han realizado las pruebas individuales de cada uno de los componentes y se ha configurado el sistema completo. En primer lugar, se proporciona una descripción general de los sensores seleccionados y se explica la importancia de monitorear los contaminantes en el aire debido a sus efectos sobre la salud humana. En segundo lugar, se llevó a cabo un monitoreo en diferentes sitios de Bogotá, en particular en las sedes de la universidad Distrital Francisco José de Caldas, donde se instalaron los concentradores junto con sus antenas. La instalación de estos dispositivos requería de energía y conectividad, por lo que se realizó una evaluación detallada de los sitios de instalación y la gestión de permisos correspondientes. Se realizó un muestreo para determinar la cobertura de la red, y se presentaron gráficos y mapas que ilustran la distribución de los sensores en la ciudad.

Posteriormente, se llevaron a cabo análisis de los datos obtenidos por los sensores y se compararon con los valores mínimos de contaminantes. Finalmente, se comparó el modelo implementado con otros modelos ya certificados de fabricantes como el sistema \gls{lwan} Multitech e ISMT, y se realizaron pruebas de interoperabilidad que resultaron exitosas. Este capítulo concluye con una evaluación detallada de los resultados obtenidos y una discusión de las implicaciones de los hallazgos para la gestión de la calidad del aire en la ciudad, la cobertura y la interoperabilidad.

\chapter{Generalidades de la redes IoT}

A pesar de que el concepto de \gls{iot}, por sus siglas en inglés, es relativamente reciente, habiendo aparecido por primera vez entre los años 2008 y 2009 \cite{Evans2011TheEverything}, y definido por el \gls{ibsg} de Cisco como “el punto en el tiempo en el que se conectaron a internet más cosas u objetos que personas” \cite{Alcaraz2014InternetCosas}, en la actualidad el IoT tiene un gran impacto en la vida cotidiana. Cada vez se desarrollan más dispositivos con conexión a internet, de modo que no solo las personas se conectan a la red, sino que también los objetos necesitan estar conectados para funcionar correctamente o cumplir con sus propósitos

\subsection*{Sobre la evaluación de la red IoT propuesta}

El \gls{iot} tiene aplicaciones tanto medioambientales como industriales, así como en ciudades, en el hogar y a nivel personal. Las aplicaciones medioambientales incluyen sensores orientados a la protección y salud, no solo del ser humano, sino también de nuestro planeta. Un ejemplo es el monitoreo de la polución en mares y ríos, midiendo el estado del agua, su PH, salinidad, temperatura, iones, entre otros. Por otro lado, existen sistemas dedicados a la protección de animales mediante geolocalización y aplicaciones para la gestión de catástrofes. Los sensores que tradicionalmente solo podían ser leídos en sus centros de control, ahora pueden comunicarse por medio de GSM y ser fácilmente accesibles (entre ellos y con las centrales). Un ejemplo destacado es el proyecto ALARMS de la BGS (British Geological Survey, 2010), que ha logrado integrarse en un dispositivo de bajo costo con capacidad GSM y sensores miniaturizados para alertar tempranamente sobre la inestabilidad del suelo, anticipando así los deslizamientos de tierra \cite{Valle2014}.

En el ámbito industrial es donde se desarrollan más aplicaciones; algunas de las actuales incluyen la monitorización de estructuras, seguridad y análisis de datos (sistemas de big data), entre otras. Las ciudades también cuentan con aplicaciones tecnológicas basadas en IoT, como en los estacionamientos, contenedores inteligentes y el control eléctrico.

En el hogar, las aplicaciones se centran en gran medida en el control del entorno, como los sistemas de riego autónomo para plantas. Por su parte, las aplicaciones personales abarcan una amplia variedad de desarrollos, incluyendo píldoras inteligentes y trajes de monitoreo para bebés \cite{Valle2014}.

Estas nuevas tecnologías, puestas a nuestra disposición, pueden ser utilizadas en casi todos los aspectos de la vida cotidiana y en la industria, gracias a las numerosas posibilidades que ofrece el IoT para obtener, procesar y transmitir información. Cuando se destaca la versatilidad de estas tecnologías en distintos ámbitos de la vida diaria, se incluyen aplicaciones como el control de electrodomésticos o cualquier otro dispositivo en el hogar capaz de recoger y/o procesar datos, monitorear el estado de la vivienda e incluso verificar nuestra salud \cite{Alcaraz2014InternetCosas}. No obstante, las aplicaciones mencionadas anteriormente benefician principalmente a los propietarios, pero también pueden extenderse a la comunidad mediante soluciones de gran alcance como redes inteligentes, luces de la calle conectadas y autos inteligentes.

\section*{Sobre la alianza \gls{lwan}}

Los sistemas de comunicaciones nacieron como una solución a las necesidades de las personas para comunicarse de forma remota; luego de  su evolución, han trascendido a las máquinas y últimamente a los objetos de interés. La comunicación de objetos se inicio para extender las redes tipo \gls{lan} y sus derivados. Luego se vio la necesidad de comunicar las máquinas de forma remota para realizar procesos síncronos, y fue allí donde surgieron las redes de telecomunicaciones celulares para prestar este servicio, creándose protocolos estandarizados para tal fin, como el M2M. La conexión de objetos bajo este protocolo se hizo costosa y no tuvo acogida. Para bajar costos y poder conectar objetos de forma masiva, hubo iniciativas privadas como Sigfox, Kore, Weighless, etc. Estas compañías se basan en tecnologías de baja potencia de área amplia LP-WAN. El éxito fue inmediato pero solo eran iniciativas empresariales.

Fue así como se inicio una alianza de los principales sectores: fabricantes, operadores, usuarios; para crear una tecnología abierta tipo  LP-WAN, la cual se llamó \gls{lora}, donde cualquier persona o empresa pueda fabricar, implementar o prestar servicios para la conexión de objetos. Las características fundamentales de LoRa son: bajo consumo de potencia, tasa de bit baja, hasta 15 kilómetros de enlace, utilización de bandas libres ISM, entre otras. 

Como parte de los avances tecnológicos, en el área del internet de las cosas, emergen nuevas propuestas que mejoran las aplicaciones actuales, es el caso de las tecnologías \gls{lpwan} \cite{LoRa-2015}, las cuales complementan las tecnologías celulares y las de corto alcance para darle inteligencia a los objetos. Aun cuando los proveedores de telefonía celular también han iniciado el mercado en este segmento,  y en especial para la tecnología \gls{m2m}, la cual puede ser extendida a la conexión de objetos. Sin embargo, se necesitaron de varios años para lograr que la industria y los operadores se pusiera de acuerdo para crear los estándares que soportaran esta tecnología, para tal fin se creo la \textit{LoRa Alliance} para definir un estándar global para redes \gls{iot} con tecnología \gls{lpwan}, \gls{m2m} y para aplicaciones industriales o de consumo.

Este proyecto es el resultado de un proceso de investigación para las redes de internet de las cosas de área amplia con tecnología \gls{lpwan}, bajo el estándar \gls{lora}. El objeto del proyecto es evaluar la red en mención, puesto  que es un estándar nuevo y de unas prestaciones que mejoran los estándares que hay en el mercado, dado que es una alianza de cientos de empresas para la conexión de objetos con sistemas tecnológicos abiertos, para la elaboración, fabricación y prestación de servicios IoT sin restricciones.

Como punto de partida, se hizo un estudio de las técnicas utilizadas en la conformación de los estándares, y lo mas relevante es el uso eficiente que se hace de la energía dado que, según las tendencias en la conexión de objetos, el crecimiento es exponencial. Para ir en consonancia en el bajo consumo de energía, se incluyo la técnica de transmisión de datos  \gls{fhss}, por su eficiencia en el manejo del ruido, como la capa física. la capa \gls{mac} que maneja las políticas del canal y la capa de aplicación. 

La red se compone de una pasarela (gateway) a la cual se le pueden conectar los nodos que le dan inteligencia a los objetos. Para el caso de este proyecto se seleccionaron como objetos, un conjunto de sensores para monitorizar ciertas condiciones del ambiente y del aire.

El estándar LoRa se define como la capa física o la modulación inalámbrica utilizada para crear enlaces de comunicaciones de rango amplio. La técnica por desplazamiento de fase, \gls{fsk}, y modulado por ensanchamiento de espectro filtrado, \gls{css}, se utiliza para  la comunicación entre nodos y pasarelas. 

La comunicación entre los nodos y la pasarela son ensanchadas en diferentes canales y con distintas tasas de datos. La velocidad de los datos está entre los 0.3 y los 50 kbps, donde los nodos se pueden programar con distintas tasas de datos en una misma pasarela, esto es posible puesto que la red LoRa puede gestionar la velocidad de los datos y la salida para cada nodo de la red por medio de un esquema adaptativo de velocidad de datos, \gls{adr}.

Algunos parámetros se especifican por regulaciones regionales y por la clase de protocolo LoRa, el cual especifica las condiciones de los protocolos de acuerdo a su aplicación, como por ejemplo, si es \textit{duplex} o \textit{fullduplex}, o el tipo de velocidades de los datos.

Los tipos de comunicaciones LoRa se enmarcan en clase A, clase B y Clase C. actualmente esta terminada y en etapa de expansión la clase A, donde los dispositivos, pasarelas y nodos, están en producción. Las otras clases están en experimentación y para este documento se hará el trabajo con los dispositivos de la clase A.

\textbf{LoRa Clase A}: En esta clase, la comunicación es bidireccional, de tal forma que a cada nodo, cuando se comunica con el servidor (pasarela), se le permite la recepción de dos ventanas de tiempo. La transmisión de las ranuras  esta basado en las necesidades de cada nodo, donde pueden tener pequeñas variaciones aleatorias similares a los protocolos de transmisión ALOHA.   luego el nodo se coloca en modo escucha y después de un retardo corto espera en dos ventanas de tiempo para recibir comunicación del servido; si recibe información en la primera ventana el nodo termina el modo de escucha. Luego este debe esperar para trasmitir datos de acuerdo a una programación definida.

\textbf{LoRa Clase B}: Es una ampliación de la clase A, donde se permite programar las recepción de mas de dos ventanas por cada nodo. Esto permite abrir ventanas extras de tiempo para un nodo. La programación de recepción para estas ventanas extras se hace con una programación en la pasarela que permite la sincronización tipo Beacon, en la cual se envía una alerta por radiodifusión desde el servidor a los nodos que están en modo escucha. En general es una comunicación bidireccional entre la pasarela y los nodos con ranuras de tiempo de recepción programadas.

\textbf{LoRa Clase C}: Para esta clase, se tienen ventanas de recepción abiertas, casi continuamente y solo están cerradas cuando el nodo transmite. Esto contrasta con un mayor consumo de potencia con respecto a las clases anteriores.

\chapter{Metodología}
Con el objeto de cumplir con todos los propósitos de este trabajo, se realizó una investigación con un enfoque experimental.

\section{Primera fase: revisión documental}

En esta fase se recolectó la información sobre el estado del arte y la información técnica referente a los dispositivos por utilizar. La documentación minuciosa permitió la selección de dispositivos sensores adecuados para la medición de la calidad del aire y del sistema completo que permitió la adquisición de las variables, su digitalización y procesamiento, para cumplir los requerimientos de las tramas en la red LoRa. 

\subsection{Etapa 1: arquitectura de la  red LoRa}
\subsubsection{Objetos/nodos }
El objetivo de esta etapa es la recolección de datos y su transmisión. Se requiere una infraestructura para la recolección y manejo de datos. El nodo es un \emph{firmware} que recolecta los datos de los objetos colocados a la entrada. Por el modelo de red, se transmite la información bajo los protocolos LoRa, con niveles de potencia muy bajos, donde el receptor es una pasarela (\emph{gateway}) que se comunica en doble vía con cada uno de los nodos que lo acceden. Los nodos y las pasarelas tienen comunicación de doble vía, lo cual los hace robustos para cualquier tipo de actuación requerida del objeto.

\subsubsection{Pasarela/concentrador}
La información que se genera en los objetos se recibe en la pasarela, la cual a su vez se puede comunicar con una gran cantidad de objetos. Para estos propósitos, se requiere la recolección, apropiación y análisis de la arquitectura de la pila de protocolos LoRa, que requiere la pasarela para una comunicación eficiente con los objetos. 

\section{Segunda fase: diseño y evaluación de los elementos de la red LoRa }

En el siguiente paso se realizaron las pruebas de los objetos para verificar las especificaciones y el comportamiento descritos por el fabricante. 

\subsection{Etapa 1: diseño de los nodos}

A partir de la selección y clasificación de la información, junto con la metodología para la evaluación de una red IoT de área amplia, se utilizaron las herramientas de simulación acordes para modelar los nodos y el sistema de comunicaciones diseñado para este propósito. Esto permitió analizar el comportamiento de los dispositivos que hacen parte de la red IoT (LoRa) ante las diferentes condiciones previamente establecidas.

\subsection{Etapa 2: diseño del módulo de procesamiento de datos}

Un nodo adquiere información de varios sensores; en el caso del proyecto, estos son detectores de agentes contaminantes del aire de por lo menos tres variables por nodo. Cada sensor requirió bloques para efectuar el muestreo en los respectivos puertos analógicos; enseguida, se realizó el procesamiento de datos según la cantidad de mediciones requeridas, en los debidos intervalos de tiempo. Adicionalmente, se diseñaron las rutinas y subrutinas necesarias para las diferentes tareas de muestreo, conversión, almacenamiento y transmisión de los datos procesados a la red LoRa.

\subsection{Etapa 3: diseño de protocolo de comunicación nodo-pasarela }

Aunque el nodo y la pasarela tienen implícita la pila de protocolos LoRa, se requirió el diseño de la interfaz de aire, puesto que la distancia marca el tipo de modulación y por ende la capacidad de la tasa de bit. Al ser un prototipo de evaluación, el diseño se basó en los parámetros tanto del nodo como de la pasarela que se eligieron para la red IoT LoRa. También se realizó la programación del gateway para la comunicación con los distintos nodos.

\subsection{Etapa 4: diseño de la  red LoRa}

Se diseñó la red IoT LoRa para la interoperabilidad de las pasarelas con la nube, con el fin de concentrar los datos en esta y posteriormente realizar la minería de datos (objeto de la segunda parte del proyecto).

\subsection{Etapa 5: simulación}

Se realizaron varias simulaciones del módulo de procesamiento de datos, dado que no se tenía conocimiento de herramientas de simulación para las redes LoRa. Los resultados de la simulación se contrastaron con el diseño y la implementación. 

\section{Tercera fase: implementación y pruebas}
Para la implementación de la red IoT LoRa, se partió de los nodos. Esto requirió tres módulos: 

\begin{itemize}
\item La implementación del módulo de lectura de datos de los sensores, que se realizó con microcontroladores de bajo costo con capacidad de manejo de cuatro o más variables.

\item La implementación del módulo de procesamiento de datos tipo estándar LoRa, el cual se evaluó de acuerdo a los siguientes parámetros: que sea un sistema abierto, que sea económico y que sea escalable para darle procesamiento en paralelo para futuros proyectos donde se requiera una variedad más amplia de variables y objetos.

\item El tercer módulo que se empleó fue el de transmisión, el cual ya tenía inmersa la pila de protocolos LoRa. 
\end{itemize}

Finalmente, y de acuerdo con el diseño, se puso a prueba la antena para la radiación de la señal hacia la pasarela.

\subsection{Etapa 1: implementación PCB}
Se diseñaron e implementaron los circuitos impresos (PCB) necesarios para los módulos del nodo; en concreto, los módulos de datos de los sensores, los cuales proveen los datos, y el módulo de procesamiento de datos con la interfaz para el nodo LoRa. También se implementó el modelo de \emph{software} que se requirió para acondicionar los datos al nodo LoRa.

\subsection{Etapa 2: pruebas del nodo}
Se conectaron los nodos a los módulos de lectura y procesamiento de datos, previa calibración de los sensores empleados. Enseguida, se hizo la transmisión de estos a través del módulo de transmisión inalámbrica.

\subsection[Etapa 3: pruebas de interoperabilidad entre pasarelas de la red IOT LoRa]{Etapa 3: pruebas de interoperabilidad\\ entre pasarelas de la red IOT LoRa}

Las pruebas se realizaron para que las pasarelas tengan comunicación en doble vía. Esto solo se hizo como parte de la evaluación y previendo próximos proyectos. 

\subsection{Etapa 4: pruebas de campo y análisis de datos}

Como el objetivo fue implementar la red IoT propuesta, se realizaron pruebas que permitieron la recolección de datos de los nodos, los cuales fueron recibidos en las respectivas pasarelas y luego concentrados en la nube para su posterior tratamiento y muestra de dicha información. Se midieron las distancias máximas que permitió la comunicación entre nodo y pasarela, en línea de vista, y para diferentes tipos de obstáculos, e inclusive para comunicación en condiciones extremas como los sótanos o subterráneos.

Las pruebas se realizaron en dos sedes de la Universidad Distrital, donde las distancias son de varios kilómetros y es óptimo para colocar el conjunto de sensores que midieron algunas de las variables contaminantes del aire en la ciudad de Bogotá.

\section[Cuarta fase: documentación y divulgación del proyecto]{Cuarta fase: documentación\\ y divulgación del proyecto}

Finalmente, se recopilaron y analizaron los resultados obtenidos, con el propósito de formular recomendaciones y conclusiones, las cuales están presentes en este documento, junto con el resto de las experiencias obtenidas. Todo lo anterior se reporta en el informe final. La divulgación del proyecto se realizó en la medida que se fueron obteniendo resultados en las diferentes etapas. Este proyecto fue presentado en un congreso internacional y adicionalmente fue publicado en una revista indexada. Por último, esperamos que sirva como punto de partida para la implementación de redes IoT en modelos de \emph{smart cities}, \emph{smart agriculture}, \emph{M2M} y redes de rango amplio para dar inteligencia a los objetos.

\newpage\thispagestyle{empty}

\chapter{Fundamentos LoRaWAN}

LoRaWAN es una \gls{wan} descrita en un modelo de tres capas: la de aplicación, la de control de acceso al medio (\gls{mac}) y la capa física o de transmisión, según se muestra en la figura \ref{f:m_protocolos}.

\begin{figure}[hb]
\centering
\caption{Capas de la red LoRaWAN}
\includegraphics[width =\textwidth]{capitulos/figuras/protocolos}
{\fuente{elaboración propia}}
\label{f:m_protocolos}
\end{figure}

\section{Capa Física}
En LoRaWAN, la capa física incluye la interfaz de radio, donde se diferencian las bandas \gls{ism} según la región, así como las técnicas de modulación \gls{fhss} y \gls{fsk}. En \gls{fhss}, el espectro se genera mediante una señal \gls{css} con variaciones continuas en la frecuencia. Esta técnica tiene la ventaja de que las compensaciones en tiempo y frecuencia son equivalentes.

Los sistemas \gls{fsk} y LoRa cuentan con modulación de envolvente constante, esquemas de bajo costo y baja potencia, con estados de alta eficiencia que pueden reutilizarse sin modificaciones. Gracias a la ganancia en procesamiento asociada a LoRa, la potencia de salida del transmisor puede reducirse y mantenerse comparable a la de un enlace convencional \gls{fsk}, al mismo tiempo que conserva o mejora las prestaciones del enlace.

El alto \gls{bt}, superior a 1, y su asincronismo natural hacen que la señal LoRa sea inmune a los mecanismos de interferencia tanto dentro como fuera de la banda. Además, proporciona excelente inmunidad frente a las interferencias de impulsos de amplitud (AM), ya que el periodo de símbolo de LoRa es más largo que las ráfagas de corta duración utilizadas en los saltos rápidos del sistema \gls{fhss}. Los valores típicos de selectividad fuera de canal son de 90~dB y el rechazo en el canal es superior a 20~dB. Los valores típicos para la modulación FSK son de 50~dB, con un rechazo de -6~dB.

El sistema también es resistente al efecto Doppler, el cual provoca pequeños desplazamientos en los pulsos LoRa, haciéndolos despreciables en el eje temporal de la señal en banda base.

Por estas razones, el sistema basado en LoRa es ideal para enlaces de comunicaciones de datos móviles, tales como sistemas de monitorización, aplicaciones de transporte, infraestructura, seguimiento, entre otros.

\section{Capa MAC} \label{secc:clases}

La capa \gls{mac} se compone de las subcapas LoRaWAN esclavo, \gls{hal} y \gls{spi}, y es la responsable de definir la clase del estándar que se implementará, pudiendo ser clase A, B o C. Este apartado aborda las especificaciones de la clase A de LoRaWAN. Asimismo, la capa \gls{mac} establece elementos clave para las redes, como la encriptación y la seguridad para \gls{m2m}, la optimización de la tasa de bits adaptativa, la calidad de servicio y otras aplicaciones avanzadas.

\subsection{LoRaWAN Clase A}

Los diferentes tipos de nodos se pueden implementar como clase A, con las capas física, MAC y de aplicación que constituyen el estándar. La definición de la clase se realiza en la capa \gls{mac}.

Una vez que el nodo transmite un mensaje, se abren dos ventanas de recepción de corta duración. La primera ventana, RX1, se activa tomando como referencia el final de la transmisión. 

\begin{figure*} 
    \centering
    \begin{tikzpicture}
        \setlength{\arrayrulewidth}{0.2mm}
        \node[draw] at (1.7, 1)   (a) {\Large{Transmisión}};
        \node[draw] at (7, 1)   (b) {\Large{RX1}};
        \node[draw] at (10, 1)   (b) {\Large{RX2}};
        \draw [black, line width=0.05cm] (0cm, 0cm) -- (3.3cm, 0cm); 
        \foreach \x in {0, 3.3} 
        \draw (\x cm, 5pt) -- (\x cm, - 5pt); 
        \draw (1.7cm, 0cm) node[below=5pt] {$t_{Tx}$}; 
        
        \draw[black, line width=0.05cm] (3.4cm, 0cm) -- (6.3cm, 0cm);
        \foreach \x in {3.4, 6.3} 
        \draw (\x cm, 5pt) -- (\x cm, - 5pt); 
        \draw (4.8cm, 0cm) node[below=5pt] {$t_{R1}$};
        
        \draw [black, line width=0.05cm] (6.4cm, 0cm) -- (9.2cm, 0cm);
        \foreach \x in {6.4, 9.2} 
        \draw (\x cm, 5pt) -- (\x cm, - 5pt);
        \draw (7.5cm, 0cm) node[below=5pt] {$t_{R2}$};
    \end{tikzpicture} 
\caption{Tiempos de comunicación nodo-pasarela-nodo}
\label{f:retardo}
\end{figure*}

Esta \textbf{primera ventana} usa la frecuencia y la tasa de datos relacionadas con las características de la transmisión. La ventana RX1 se abre tras un retardo, $t_{R1}$, de aproximadamente \SI{20}{\us}, como se ilustra en la figura \ref{f:retardo}. Este retardo varía de acuerdo con las bandas regionales especificadas en los parámetros regionales de \gls{lwan}. Por defecto, la tasa de datos en RX1 es igual a la de la transmisión.

La \textbf{segunda ventana}, RX2, tiene parámetros de frecuencia y tasa de datos fijos, aunque pueden ser modificados mediante comandos \gls{mac}. El retardo $t_{R2}$ es de aproximadamente \SI{20}{\us}, como se muestra en la Figura\ref{f:retardo}. Los valores predeterminados de estos parámetros están definidos por las bandas regionales de los estándares \gls{lwan}.

La duración de la ventana de recepción debe ser al menos suficiente para que el transceptor del nodo pueda detectar el preámbulo del mensaje entrante.

\subsection{LoRaWAN Clase B}

Los nodos y la pasarela tienen comunicación bidireccional con ranuras de recepción programadas. Los nodos de clase B, además, tienen las funcionalidades de la clase A. Para alcanzar la sincronización de la red, los nodos reciben un tiempo de sincronización tipo Beacon desde la pasarela, esto permite al servidor saber cuándo un nodo está escuchando. Los nodos que hablen en la ventana de tiempo programada entran a la pasarela de acuerdo a procesos aleatorios con prioridad en servicios a los nodos en mención. 

\subsection{LoRaWAN Clase C}

Los nodos y la pasarela tienen comunicación bidireccional con ranuras de recepción máximas. Proporciona ventanas de recepción casi continuamente, solo se cierra cuando el nodo está transmitiendo. Este modo se utiliza cuando hay requerimiento de grandes cantidades de datos entre el nodo y la pasarela. 

\section{Tramas de transmisión y recepción}

Las tramas de la capa física se distinguen entre las de transmisión (\textit{uplink}) y las de recepción (\textit{downlink}), desde el punto de vista del nodo.

\begin{figure}
\caption{Estructura de tramas de \textit{uplink} y \textit{downlink}}
\centering
    \resizebox*{.8\textwidth}{2cm}{
    \begin{tabular}{|c|c|c|c|c|}  \hline
        \multicolumn{5}{|c|}{{\textbf{Estructura de transmisión} }}\\ \hline
         Preamble & PHDR & PHDR-CRC & PHYPayload & CRC \\ \hline \hline
         \multicolumn{5}{|c|}{{\textbf{Estructura de recepción} }}\\ \hline
         Preamble & PHDR & PHDR-CRC & PHYPayload& -- \\ \hline
    \end{tabular}}
\label{f:retardo1}
\end{figure}

\enlargethispage{-\baselineskip}

El nodo transmite la trama a la pasarela, figura \ref{f:retardo1}, y de ahí al servidor de red. La trama de la capa física de LoRa está compuesta por el preámbulo, el \gls{phdr}, el \gls{pc}, la carga útil (\textit{payload}) y el \gls{crc}. 

La trama de recepción la envía la pasarela a un solo nodo y es retransmitida por una única pasarela. Los mensajes de la capa física también se denominan modo explícito de paquetes de radio. Los mensajes de recepción tienen la misma estructura del mensaje de transmisión, a excepción del \gls{crc}, el cual está ausente, como se muestra en la parte baja de la figura \ref{f:retardo}.

La estructura de las tramas LoRa de la capa física contiene la información de las capas superiores (figura \ref{f:m_lora}), tal como lo es la capa \gls{mac}. El asterisco hace referencia a que el \gls{crc} solo está presente en la trama de transmisión.

\begin{figure}
    \caption{Formato del mensaje LoRa}
    \centering
    \includegraphics[width =\textwidth]{capitulos/figuras/f2}
    \label{f:m_lora}
\end{figure}

En la \gls{pp}, se encuentra el contenido de los diferentes tipos de tramas \gls{mac}, las cuales tienen en común el encabezado \gls{mhdr} y el \gls{mic}. Los tipos de tramas \gls{mac} corresponden a los de la \gls{mp}, solicitud de acceso (\textit{Join-Request}) y respuesta a la solicitud (\textit{Join-Response}). 

La trama \gls{mp} contiene la información de los nodos. Los campos de esta trama corresponden al \gls{fhdr}, \gls{fp} y a la \gls{fpa}. La \gls{fhdr} contiene los campos de \gls{da}, de \gls{fc}, de \gls{fcn} y las \gls{fo}.

\section{Características de LoRaWAN}

\textbf{Banda de frecuencias y canalización}. La banda EU-868 del estándar LoRa tiene un rango de \SI{863}{\MHz} a \SI{870}{\MHz} y también está definida en el espectro de radio para la banda \gls{ism} normatizada por \gls{etsi}, en el estándar EN300.220. Los canales de la red pueden ser atribuidos libremente por el operador. La pasarela posee ocho canales en paralelo, IF0 al IF7, como medio de comunicación con el nodo. El ancho de banda es fijo, 125~kHz, con tasa de bit variable de acuerdo al \gls{sf}.

\begin{figure}
    \caption{Canalización banda ISM868}
    \centering
    \includegraphics[width =\textwidth]{capitulos/figuras/canales}
    \label{f:m_canales}
\end{figure}

Además tiene dos canales adicionales, uno LoRa con ancho de banda fijo de 250 kHz con una tasa de bit de 11~kbps como el canal \textit{back-haul} de la red LoRaWAN IF1 y otro canal de \gls{gfsk}, el cual tiene una tasa de datos y un ancho de banda ajustables para una velocidad máxima de 50 kbit/s \cite{LoRa-2015}.

\textbf{Factor de ensanchado \gls{sf}, tasa de bit y potencia}. Permite un manejo eficiente del espectro y de la capacidad de la red LoRa. Hay seis factores de ensanchado \gls{sf}, que van del SF7 al SF12, los cuales están relacionados con la potencia y la tasa de datos. La tasa de datos va desde 250~bps, para \gls{sf}12 con una sensibilidad de recepción de -137~dBm, hasta 5.5~kbps, para \gls{sf}7 con una sensibilidad de -123~dBm. La potencia máxima de transmisión está, de acuerdo a \gls{etsi}, en 14~dBm. LoRa utiliza el ajuste dinámico de la potencia de emisión y la \gls{adr}, de acuerdo a la distancia que separa el nodo de la pasarela, de esta manera se maneja la eficiencia de la energía.

\section{Seguridad LoRaWAN}

La seguridad en LoRaWAN es uno de los activos más importantes en aplicaciones IoT, industriales, \gls{m2m}, ciudades inteligentes, etc. La seguridad cubre múltiples propiedades y en especial las criptográficas, las cuales merecen un estudio más profundo.  

La autenticación es bidireccional entre la red y el nodo para el flujo de los datos, una vez activado el nodo. Esto asegura que los nodos se unan y se autentiquen a la respectiva red. 
Los mensajes de aplicación de la capa \gls{mac} de LoRaWAN tienen autenticación de origen, integridad protegida, reproducción protegida y encriptación. Lo anterior unido a la autenticación mutua garantiza que el tráfico en la red no sea alterado y que proviene de un nodo legítimo de la red. Esto hace que la información no sea comprensible para los intrusos y tampoco permite ser capturada o alterada por los piratas. La tecnología LoRa es una de las pocas redes IoT que implementa encriptación de extremo a extremo.  
En las redes celulares el tráfico es encriptado en la interfaz de aire, pero es transportado como texto plano en el núcleo de la red del operador, lo cual puede ser vulnerado de forma selectiva. Para este problema se creó una etapa adicional para la seguridad. LoRaWAN no provee esta capa adicional ya que se aumentaría el consumo de potencia, la complejidad y los costos.

Potentes algoritmos de seguridad están implementados en la red LoRa como lo es el \gls{aes}, el cual ha sido adoptado para la seguridad del tráfico entre los nodos y la red en conjunto a los diferentes modos de operación como lo son: \gls{cmac} para proteger la integridad de la información y el \gls{crt} para la encriptación. Los nodos son personalizados con una clave \gls{aes} de 128~bit (AppKey) y un único identificador global (EUI-64-based DevEUI), figura \ref{f:seg1}. 

\begin{figure}
    \caption{Seguridad en la trama LoRaWAN}
    \centering
    \includegraphics[width =\textwidth]{capitulos/figuras/seguridad1}
    \label{f:seg1}
\end{figure}

La activación al aire (\textit{over-the-air}) provee el reconocimiento de ambos, el nodo y la red, por medio de la llave \textit{AppKey} el nodo se puede unir a la red, figura \ref{f:seg2}. Dos llaves de sesión que corresponden: una a la integridad del tráfico y la otra a la encriptación de los comandos y la carga útil de la aplicación de la capa \gls{mac} de \gls{lwan}  llamada \textit{NwkSKey}, y otra para la encriptación de extremo a extremo de la carga útil de la aplicación denominada \textit{AppSKey}. 

\begin{figure}
    \caption{Seguridad LoRaWAN}
    \centering
    \includegraphics[width =\textwidth]{capitulos/figuras/seguridad2}
    \label{f:seg2}
\end{figure}

La llave \textit{NwkSKey} es distribuida por la red LoRaWAN en su orden para probar y verificar la autenticidad y la integridad de los paquetes. La llave \textit{AppSKey} la provee el servidor de la aplicación para encriptar y desencriptar la carga útil de la aplicación. Las llaves, \textit{AppSKey} y \textit{AppKey}, están ocultas al operador de la red, por tal razón no se puede desencriptar la información del usuario.

\section{Capa de Aplicación}

Una \gls{wsn} es un sistema de comunicación conformado por varios nodos sensores, donde cada uno de ellos está dotado de los elementos apropiados para la detección de magnitudes físicas asociadas a procesos industriales, a la calidad del aire, al medio ambiente, o a otros eventos en la naturaleza.

La \gls{wsn} se ha hecho realidad por la convergencia de los desarrollos tecnológicos en los \gls{mems}, las comunicaciones inalámbricas y la electrónica digital.

La tecnología de miniaturización de los nodos para \gls{wsn}, basada en \gls{mems}, ha hecho un progreso notable en los últimos años.

La tecnología de \gls{mems} consiste en combinar la microelectrónica, el micromecanizado y el empaquetado en alta densidad.

Las estructuras en niveles 2D y 3D pueden ser producidas y soportadas por tecnologías de microelectrónica y micromecanizado, para pasar a formar parte de los sistemas sensores en miniatura, los cuales consumen baja potencia y llevan incluido directamente su sistema de acondicionamiento.

En este momento existe una gran variedad de productos \gls{mems} en el mercado, los cuales permiten la detección de una variedad de magnitudes en áreas de la física, la biología, la biomasa, entre otras. Entre dichas magnitudes se destacan los sensores para desplazamiento, velocidad, aceleración, presión, tensión, sonido, luz, variables eléctricas, variables termodinámicas, variables biológicas, y muchas más.

La tecnología de los \gls{mems} ha llegado, en este momento, a desarrollar nodos sensores de dimensiones del orden de $\SI{2}{mm}\times \SI{2}{mm}\times \SI{1}{mm}$.

La combinación de \gls{iot} y \gls{wsn} permite un desarrollo revolucionario de los sistemas de comunicación con mejoras en la confiabilidad, flexibilidad, cubrimiento y eficiencia.

La \gls{wsn} está configurada por una red de nodos que cooperativamente sensan y controlan un medio, lo que permite la interacción entre las personas o los sistemas de cómputo y el entorno \cite{Akyildiz-2002}.

La actividad de detección, procesamiento y comunicación está limitada por la cantidad de energía disponible, lo cual exige un diseño integral que conlleva procesamiento de señales, control de acceso al medio y protocolos de comunicación.

Las aplicaciones de la tecnología \gls{wsn} permiten el desarrollo de sistemas inteligentes en áreas tales como manejo de aguas, transporte, hogar, industria, entre otras y siempre con mejoras importantes sobre las tecnologías tradicionales.

Agregado a lo anterior, el mundo moderno de las \gls{wsn} emigra hacia una nueva era con su integración al \gls{iot}, donde habrá una gran variedad de temas que deberán ser normados y aceptados en un corto tiempo. Algunos de estos temas más apremiantes son la propiedad, uso y legalidad de los datos que se generan, directa o indirectamente, gracias a dichos sistemas.

En su estructura fundamental, una \gls{wsn} está constituida por:
\begin{itemize}
\item Los nodos sensores dispuestos para detección de las diferentes magnitudes y para la comunicación.
\item Los nodos actuadores, los cuales ejecutan funciones de control comandadas de acuerdo al resultado de una medición previa.
\item La pasarela o gateway encargada de la conexión a internet.
\item Los clientes. 
\end{itemize}

Los nodos transmiten y reciben el estado de cada uno de ellos entre sí para reconocerse como existentes en la red. Los nodos configuran una red (estrella, o malla, o cualquier otra) con cubrimiento de distancias cortas con un máximo de 1000~metros y utilizando la técnica multi-salto. 

\newpage\thispagestyle{empty}
 
\chapter{Nodos LoRa}

En términos de redes de telecomunicaciones móviles, a los dispositivos que consumen la red inalámbrica se los denomina comúnmente como \emph{nodos}, y es así que en \gls{lwan} se hablará del \emph{dispositivo \gls{lora}}, también conocido como \emph{mote}. De igual manera, en el modelo cliente/servidor, el nodo se comunicará con el servidor \gls{lwan}, evitando la comunicación directa con otros dispositivos de la red. El servidor \gls{lwan} podrá enviarle de vuelta al nodo mensajes almacenados en el servidor \gls{lora}, siempre dependiendo de la clase de trabajo en que se encuentre el nodo (sección \ref{secc:clases}). Para efectos prácticos, se hablará del nodo funcionando en clase A, en la que el nodo siempre inicia la comunicación. Las demás clases conservan el mismo modelo cliente/servidor, por lo que no se abordarán en detalle.

La composición del nodo se basa típicamente en una arquitectura de \emph{hardware} de referencia (figura \ref{fig:arqgenhw}), la cual se puede abordar según la funcionalidad de cada uno de los bloques o según la implementación disponible de la arquitectura base mencionada.

\section{Bloques de la arquitectura genérica}

En la arquitectura genérica de la figura \ref{fig:arqgenhw}, cada bloque tiene una función específica dentro del nodo:
\begin{itemize}
    \item La fuente de energía puede ser suministrada mediante un conector o una batería.
    \item La \gls{mcu}, o simplemente el \gls{uc}, maneja todas las funciones del dispositivo e implementa la pila LoRaWAN.
    \item El radio LoRa está compuesto por el transceptor \gls{lora}, el circuito de acople de la antena y la antena en sí misma.
    \item Los periféricos pueden ser sensores como acelerómetros o sensores de temperatura, o entradas y salidas tales como relés o pantallas.
\end{itemize}

\begin{figure}[H]
    \caption{Arquitectura de \emph{hardware} de un dispositivo \gls{lora} genérico}
    \label{fig:arqgenhw}
    \centering
    \includegraphics[width=0.8\textwidth]{arq-gen-hard}
    \propia
\end{figure}

\subsection{Fuente de poder}

En un \gls{uc} generalmente se observa que el consumo de energía es directamente proporcional a la velocidad de procesamiento y a la cantidad de periféricos activos, entre otros factores. 

Esto es crítico en aplicaciones que funcionan la mayor parte del tiempo, como una \gls{wsn} (capítulo \ref{cap:aplicacion}), por lo cual se hace conveniente disminuir el consumo. Una forma de lograrlo es haciendo que el nodo pase la mayor parte del tiempo en un estado ``adormecido'', esto es, un estado en que su actividad de procesamiento es mínima. Dependiendo del tipo de \gls{uc}, de la arquitectura y de los diversos estados de adormecimiento, la corriente requerida suele oscilar entre algunos microamperios y cientos de microamperios. El consumo de energía en trabajo normal es una medida útil y se hace más importante cuanto más tiempo pasa un nodo en periodos de trabajo normal, donde los consumos exigen corrientes que oscilan entre \SI{1}{mA} y \SI{10}{mA} para \gls{uc} pequeños y de bajo consumo, pero pueden variar dependiendo de los tipos de periféricos incorporados, la velocidad del reloj, el tamaño de la memoria y la localización del programa de ejecución.

La tarea primordial del nodo sensor es detectar eventos, realizar el procesamiento básico de los datos y transmitir la información. Lo anterior implica que en el consumo de energía se destacan tres aspectos: sensado, procesamiento y comunicación. La comunicación incluye transmisión y recepción, aspectos que demandan mayor consumo. \cite{Akyildiz-2002}.

La energía será usualmente un recurso limitado en el nodo sensor, lo que implica el uso de fuentes de energía alternativas tales como celdas solares o generadores eólicos, dependiendo del tipo de ambiente en el cual se halla el módulo sensor.

Otros aspectos relativos al suministro de energía están ligados al tiempo de vida de las fuentes y al acceso al nodo sensor. 

\subsection{Microcontrolador}

El \gls{uc} se puede considerar como el centro de control del nodo sensor. Casi todas las familias de \gls{uc} disponibles son diferentes: algunos pueden ser diseñados para la velocidad, y otros para la eficiencia energética. En la cuestión de elegir el \gls{uc} correcto para la aplicación particular del nodo, hay varios aspectos para tener en cuenta en el momento de diseñarlo.

A continuación, se presentan los componentes más relevantes del \gls{uc}:

\begin{itemize}
	\item Memoria:    
    la cantidad y el tipo de memoria utilizados son dos variables que están ligadas al consumo de energía y que deben ser cuidadosamente elegidas, dado que ambos parámetros implican costos por controlar. En general, a mayor cantidad de RAM, más potencia requerida para su sostenimiento; no obstante, los \gls{uc} modernos pueden incorporar varios cientos de kilobytes de memoria \emph{flash} y decenas de kilobytes de RAM y continúan utilizando paquetes estandarizados de manejo con costos bajos.
    
    Una tecnología nueva en este aspecto es el uso de FRAM, una alternativa de memoria que combina la velocidad de la memoria RAM convencional con la no volatilidad del almacenamiento de la memoria \emph{flash} y que elimina la separación entre las dos. Este tipo de memoria tiene menor consumo, cuando funciona, en comparación con la memoria \emph{flash} estándar, y elimina la necesidad de refrescamiento de la RAM para preservar su contenido.

    \item Soporte de entrada-salida: 
    el \gls{uc} está en capacidad de llevar a cabo funciones de comunicación tipo serial o paralela, temporizaciones, comparaciones, transmisión y recepción digital y analógica en diversos protocolos y velocidades, además del procesamiento de datos. Una de las limitaciones a estas funciones está en el número de pines de acceso físico a su estructura. Para dar solución a esta situación, se trata de mantener al máximo el número de funciones y se multiplexa la comunicación de  entrada-salida en los pines disponibles, seleccionando a través de \emph{software} la gran variedad de operaciones que permita el dispositivo.

    Para el diseño de un nodo sensor se requieren muchas de las funciones presentes en el \gls{uc} y una cierta cantidad de pines de comunicación con funciones particulares. Como ejemplos de funciones y pines requeridos, se tienen las interfaces \IC, \IS, canales SPI, conversores ADC y DAC, pines de interrupciones E/S, señales de tiempo y de reloj, bits de bandera, entre otros.

    \item Soporte de comunicación: 
    los nodos sensores están normalmente dispersos en áreas específicas de medición. Cada uno de estos nodos tiene la capacidad de recopilar datos y enviarlos a la pasarela y/o a los usuarios finales.

    Los protocolos utilizados por el receptor y por los nodos de sensores permiten el encaminamiento e integración de los datos con los protocolos de red de manera eficiente a través del medio inalámbrico, distribuyendo el trabajo de los nodos sensores.

	Los nodos sensores están debidamente coordinados y administrados para que realicen tareas que disminuyan el consumo total de energía; por ejemplo, informar a sus vecinos cuando su nivel de energía es bajo y por tanto no puede participar en el intercambio de mensajes.

    \item Adquisición de datos: 
    el modelo que se siga para la toma de datos caracteriza y describe la interacción entre los sensores y la aplicación.

    El modelo de adquisición de datos más apropiado para una aplicación se determina según las características que se desee dar a la aplicación; en las diferentes situaciones, los datos fluyen principalmente desde el nodo del sensor hasta la puerta de enlace; en las aplicaciones de control, los datos, adicionalmente fluyen desde la puerta de enlace hacia los nodos sensores.

	Las técnicas de monitoreo más usuales \cite{Cecilio-2014} son:
	\begin{itemize}
 		\item Muestreo periódico: para aplicaciones donde las condiciones o los procesos requieren ser monitoreados constantemente, como es el caso de la temperatura en un ambiente acondicionado o la presión en una tubería de proceso. Para este caso, los datos se adquieren en el nodo sensor y son reenviados a la puerta de enlace de forma periódica. El periodo de muestreo depende de qué tan rápido varía la magnitud monitoreada y/o qué características de control se requieren en el proceso. Si la dinámica del proceso controlado varía, los nodos sensores deben ser capaces de adaptar sus tasas de muestreo a la dinámica del proceso. Esta adaptación debe hacerse automáticamente o a través de comandos emitidos por la aplicación.
		\item Muestreo impuesto por el evento: hay muchos casos que requieren el monitoreo de variables críticas, como las alarmas de fuego, las alarmas de presión o los sensores de puertas o ventanas, etc. La detección de dichos eventos exige que el nodo sensor actúe con eficiencia en el consumo de energía y con la adecuada velocidad de respuesta.
		\item Captura-proceso-almacenamiento: en algunos casos, los datos pueden ser capturados y almacenados o incluso procesados por el nodo sensor antes de que se transmita a la puerta de enlace, en lugar de transmitir inmediatamente cada muestra a medida que se adquiere, lo cual mejora el rendimiento de la red en cuanto a eficiencia de consumo y ancho de banda.
	\end{itemize}
\end{itemize}

\subsection{Radio LoRa}

La comunicación de radio puede ser bidireccional, tanto entre los nodos como hacia las estaciones base. Adicionalmente, se debe buscar un balance entre el consumo de energía del módem \gls{lora} y la \gls{snr}. Dependiendo de este balance, habrá un impacto consecuente en el consumo de energía de este elemento del sistema, comparado con los otros, sobre todo en el momento en que el nodo requiera transmitir información.

\subsection{Periféricos}

Los sensores y actuadores son la razón de ser del nodo. Estos se distribuyen dentro del área de cobertura del nodo para supervisar y recopilar datos con el objetivo de transmitirlos a otros nodos y/o a un servidor.

\section{Topologías}

Dependiendo de los requisitos de diseño y producción, se tienen varias opciones para construir un dispositivo \gls{lora}. La elección de la arquitectura objetivo se hace con base en criterios como las cantidades de producción esperadas, las habilidades disponibles en ingeniería de RF y la línea de tiempo de desarrollo disponible para completar el proyecto.

Con el fin de seleccionar la arquitectura óptima, es preferible partir de un kit de desarrollo \gls{lora} y crear un bosquejo.

Por otro lado, en la figura \ref{fig:arqgensw} se muestra una arquitectura genérica del \emph{software} en el dispositivo.

\begin{figure}[t]
    \centering
    \caption{Arquitectura genérica del \emph{software} del nodo}
    \label{fig:arqgensw}
    \includegraphics[width=\textwidth]{arq-gen-soft}
    \propia
\end{figure}

La capa de \emph{drivers} provee la adaptación de \emph{hardware} e implementa todos los \emph{drivers} y \emph{buffers} para controlar los dispositivos periféricos. Esta permite abstraer el \emph{hardware} como funciones simples expuestas al \emph{middleware}.

El \emph{middleware} implementa las librerías de protocolos de comunicación (\gls{lwan}, 6LowPAN, etc.) y los controladores complejos, como el control de pantallas y de GPS.

La capa de aplicación contiene todas las aplicaciones funcionales en las que se implementa el comportamiento del dispositivo y sus funcionalidades.

\subsection{Diseño basado en un chipset LoRa}

Este diseño sigue la arquitectura genérica de la figura \ref{fig:arqgenhw}, con el radio \gls{lora} basado en un transceptor \gls{lora} de Semtech y la antena, más su circuitería asociada.

\begin{figure}[t]
    \centering
    \caption{Arquitectura de dispositivo basado en un chipset LoRa}
    \label{fig:loraWchipset}
    \includegraphics[width=\textwidth]{design-chipset}
    \propia
\end{figure}

En esta arquitectura, ilustrada en la figura \ref{fig:loraWchipset}, el desarrollador del dispositivo deberá tener cuidado de todo el diseño RF, incluyendo las restricciones de enrutamiento en la PCB, el acoplamiento y sintonización de la antena y los problemas de emisión e inmunidad.

La totalidad de la pila \gls{lwan} será llevada por el MCU principal e implementada por el desarrollador o fabricante del dispositivo.

\subsection{Diseño basado en un módulo LoRa}

El módulo \gls{lora} es un componente que contiene el \gls{mcu} y el radio \gls{lora}. Se tiene la disponibilidad de programar el \emph{software} de su \gls{mcu} para ejecutar la aplicación y la pila \gls{lwan}.

El principal beneficio de este enfoque de diseño es que todo el desarrollo del \emph{hardware} de RF es implementado por el fabricante del módulo.

La mayor parte de la sintonización y acoplamiento de la antena es realizada dentro del módulo. El fabricante del módulo brinda un diseño de referencia para conectar o rediseñar la antena.

\begin{figure}[t]
    \centering
    \caption{Arquitectura de dispositivo basado en un módulo LoRa}
    \label{fig:loraWmodule}
    \includegraphics[width=\textwidth]{design-module}
    \propia
\end{figure}

En comparación con la arquitectura anterior, dado que se dispone del MCU, en esta arquitectura (figura~\ref{fig:loraWmodule}), el fabricante del dispositivo se encarga solamente de la implementación de la pila \gls{lwan}. Dependiendo del proveedor del módulo, la pila \gls{lwan} se puede entregar como una librería.

\subsection{Diseño basado en un RF-MCU}

Un RF-MCU es un \gls{soc} que incluye un \gls{mcu} y un transceptor LoRa en el mismo circuito integrado. Los requerimientos RF son los mismos del diseño RF anterior, pero esta arquitectura ofrece mayor ganancia en el tamaño total del dispositivo. Desde el punto de vista del desarrollo, este diseño es similar al del módulo.

\subsection{Diseño basado en un módem LoRa}

Un módem LoRa es un componente que contiene todos los elementos de radio: pila y circuitería RF (figura \ref{fig:loraModemArch}). Algunos módems pueden también incluir una antena integrada, o el fabricante da una referencia de diseño para diseñar o conectar una antena externa.

\begin{figure}[t]
    \centering
    \caption{Arquitectura de dispositivo basado en módem LoRa}
    \label{fig:loraModemArch}
    \includegraphics[width=0.9\textwidth]{design-modem}
    \propia
\end{figure}

El módem integra la pila \gls{lwan} entera y actúa como esclavo, de tal manera que requiere de un \emph{host} para manejarlo. La comunicación usualmente se realiza mediante comandos AT para configurar o enviar mensajes.

La interfaz de \emph{hardware} puede ser usualmente UART, USB o SPI.

\subsection{Dispositivos con un módem LoRa externo}

Un módem LoRa también puede ser un dispositivo independiente externo (figura~\ref{fig:loraExternal}), conectado a un dispositivo existente mediante una conexión USB, por ejemplo.

En ese caso, el dispositivo existente debe tener una conexión externa disponible y también debe ser capaz de energizar el módem LoRa externo.

El módem integra la pila \gls{lwan} entera y actúa como un esclavo, de tal modo que requiere de un controlador externo para manejarlo. La comunicación usualmente se realiza mediante comandos AT para configurarlo o enviar mensajes.

\begin{figure}[H]
    \centering
    \caption{Arquitectura de dispositivo basado en módem LoRa externo}
    \label{fig:loraExternal}
    \includegraphics[width=0.9\textwidth]{design-external}
    \propia
\end{figure}

\section{Interfaces digitales}

Los sensores disponibles en el mercado pueden requerir polarización y acondicionamiento, como pueden ya tener integrados estos circuitos en el mismo encapsulado e inclusive presentar una salida con alguna interfaz digital comercial.

Debido a la variedad tanto de interfaces digitales comerciales como de distintas configuraciones de polarización para los sensores que requieren circuitos de polarización y/o acondicionamiento, se plantea una interfaz que permita conectar la mayor cantidad posible de componentes necesarios entre el \gls{uc} conectado al nodo LoRaWAN y el sensor que se vaya a conectar al nodo.

Los protocolos de comunicación electrónica más comunes entre componentes de sistemas embebidos suelen ser mediante buses serie de varios tipos. Se hablará solamente de las líneas dedicadas exclusivamente al protocolo en mención, excluyendo, así, las de polarización y energización de los dispositivos.

\subsection{Inter-Integrated Circuits (I2C)}

El bus \gls{i2c} se desarrolló en 1982 por Phillips (actualmente, NXP). Se utiliza principalmente en la comunicación de dispositivos distribuidos en un mismo sistema embebido, por ejemplo, entre un \gls{uc} y circuitos periféricos integrados. Es un bus maestro-esclavo, en el que la transferencia de datos es siempre inicializada por un maestro y el esclavo reacciona ante esta (figura~\ref{fig:i2c}).

\begin{quote}\small
    Es posible tener varios maestros mediante un modo multimaestro, en el que se pueden comunicar dos maestros entre sí, de modo que uno de ellos trabaja como esclavo. El arbitraje (control de acceso al bus) se rige por las especificaciones, de modo que los maestros pueden ir turnándose. (Aguilar Castillejos, 2019, p. 1)
\end{quote}
    
Requiere solamente dos líneas, para reloj y datos. Ambas se encuentran en un estado lógico alto débil (debido a resistencias \emph{pull up} en cada línea), y por lo mismo quienes usen el bus han de tener salidas tipo colector/drenador abierto. Quien inicia una transmisión (tomando el control de las dos líneas) es siempre un dispositivo maestro, y puede iniciarla con uno de los demás dispositivos en el bus, que se comportarán como esclavos dependiendo del direccionamiento.

\begin{figure}[H]
    \centering
    \caption{Esquema \IC}\label{fig:i2c}
    
    \includegraphics[width=\textwidth]{I2C}

    \nota{realmente, la polarización (líneas VCC y GND) puede ser distinta para cada dispositivo, pero las líneas del bus, SDA y SCL, son comunes a todos eléctricamente.}
    
    \vspace{.3\baselineskip}
    \fuente{\url{http://www.prometec.net/bus-i2c/}}
\end{figure}

\subsection{SPI}

También en modelo maestro-esclavo, con comunicación iniciada por el dispositivo maestro, este bus permite comunicación \emph{full duplex} mediante las líneas MOSI (\emph{master out, slave in}) y MISO (\emph{master in, slave out}), con reloj y habilitación del esclavo establecidos por el maestro, mediante las líneas CLK y SS (\emph{slave select}) (figura~\ref{fig:spi}).

Este bus permite velocidades de transmisión más rápidas, tanto porque las líneas son activas (requieren salidas tipo \emph{push-pull}) como porque el maestro puede enviar y recibir al tiempo, a costa de requerir más pines que \IC\ o UART.

\begin{figure}
    \centering
    \caption{Esquema SPI}\label{fig:spi}
    \includegraphics{SPI}
    \nota{todos los esclavos comparten y se conectan a las mismas líneas del bus: SCLK, MOSI y MISO, pero solo son activos bajo su señal SS.}

    \vspace{.3\baselineskip}
    \fuente{Cburnett (2006).}
\end{figure}

\subsection{UART}

El transmisor-receptor asíncrono universal\foreignlanguage{english}{(\emph{universal asynchronous receiver-transmitter})} es similar a SPI dado que permite comunicación \emph{full duplex}, mas no utiliza ni las líneas de reloj (dado que cada bit tiene un tiempo determinado y en mutuo acuerdo) ni de habilitación (cualquiera de los dos puede iniciar una transmisión). Como mínimo, cada dispositivo que soporte UART presenta las líneas Rx y Tx; para comunicar dos dispositivos UART basta con intercambiar las líneas (figura~\ref{fig:uart}).

\begin{figure}[t]
    \centering
    \caption{Esquema UART}
    \label{fig:uart}
    \includegraphics[width=0.6\textwidth]{uart}
    \nota{las líneas de los dispositivos se conectan intercambiadas, emparejando siempre un receptor con un transmisor.}
    
    \vspace{.3\baselineskip}
    \fuente{\url{https://www.mikroe.com/blog/uart-serial-communication}}
\end{figure}

\section{Interfaz planteada}

La interfaz abarca dos entornos: uno de \emph{hardware}, relacionado con la conexión física, y otro de \emph{software}, relacionado con una capa de abstracción entre los distintos métodos de obtención de las variables de los sensores, ya sea mediante protocolos digitales o mediante una cadena de medida, con el fin de permitir su transmisión sobre una aplicación en la red \gls{lwan}.

Con ayuda de una aplicación \gls{lwan}, esta interfaz permitirá la interconexión de dos circuitos: uno de estos será el núcleo Pylates, conformado por un \gls{mcu} y un módem \gls{lwan}. El otro será la placa de aplicación, con todos los circuitos adicionales requeridos en el nodo.

\subsection{Interfaz de \emph{hardware}: conector y funcionalidades de pines}

Ha de ser posible que el \gls{mcu} acceda a las variables de medida de cada sensor, bien mediante interfaces digitales (SPI, \IC, UART), bien con la implementación de una cadena de medida (OPAMPS+REF+ADC+DAC) para aquellos en que sea necesario. Se estima el peor de los escenarios, en el que se requiere la mayor cantidad de pines: cuatro sensores conectados al \gls{mcu}, tres mediante las interfaces digitales mencionadas y uno con una cadena de medida asistida por los periféricos analógicos del \gls{mcu}. Otra alternativa consiste en considerar solo un tipo de sensor, siendo en este caso conveniente permitir a los sensores que requieren cadena de medida acceder a la mayor cantidad de periféricos analógicos disponibles del \gls{mcu} (figura~\ref{fig:con_main}).

\begin{figure}[H]
    \centering
    \caption{Conector base propuesto para el nodo}
    \label{fig:con_main}
    \includegraphics[width=0.25\textwidth]{conector_principal}

    \nota{Los seis pines inferiores son la versión reducida; los demás hacen parte de la versión extendida.}

    \vspace{.3\baselineskip}
    \propia
\end{figure}

Dado que ambos conectores presentan pines en común, se puede escoger cualquiera de los dos según la complejidad del circuito esclavo por controlar.

\begin{table}[H]
\centering
\caption{Disposición sugerida de las funcionalidades de los pines}
\footnotesize
\label{tab:conector_pines}
\begin{tabular}{llllll}
\toprule
\textbf{Pin} & \textbf{Tipo} & \textbf{Función por defecto} & \textbf{Pin} & \textbf{Tipo} & \textbf{Función por defecto}          \\ \midrule
1   & PWR  & VCC, 2.8 V - 3.6 V  & 2   & PWR  & GND                                                                      \\ \midrule
3   & I/O  & VREFH                                                           & 4   & I/O  & VREFL                                                                    \\ \midrule
5   & I/O  & --Bloqueado--                                                   & 6   & I/O  & \parbox{3cm}{(CLK*/SCK)Op2,\\salida}        \\ \midrule
7   & I/O  & Op1, salida                                                     & 8   & I/O  & \parbox{3cm}{(MOSI*/Tx*)Op2,\\inversora}    \\ \midrule
9   & I/O  & Op1, inversora      & 10  & I/O  & \parbox{3cm}{(MISO*/Rx*)Op2,\\no inversora} \\ \midrule
11  & I/O  & Op1, no inversora   & 12  & I/O  & SS*(SDT)                                                                 \\ \midrule
13  & I/O  & Tx*                                                             & 14  & I/O  & MOSI*                                                                    \\ \midrule
15  & I/O  & Rx*                                                             & 16  & I/O  & MISO*                                                                    \\ \midrule
17  & I/O  & SDT*                                                            & 18  & I/O  & SCK/CLK*                                                                 \\ \bottomrule
\end{tabular}

\pieimagensimple{* Pines para periféricos digitales, que pueden ser asignados a otros pines, y entre paréntesis para el conector base.}
\end{table}

%%%%%%%%%%%%%%%%%%%%%%%%%%%%%%%%%%%%%%%%%%%
\subsection[Interfaz de \emph{software}: configuración mediante ID de núcleo y placa]{Interfaz de \emph{software}: configuración\\ mediante ID de núcleo y placa}

En aras de que el usuario final realice la mayor cantidad de configuración desde una interfaz gráfica y con un mínimo esfuerzo, se plantea una aplicación en que se ingresan dos identificadores: uno del núcleo y otro del circuito de aplicación.

\subsubsection{Núcleo}

Cada núcleo, conformado por un microcontrolador y un módem LoRaWAN, estará asociado a un identificador universal del núcleo (CoreID). Este identificador permite su indexación en una base de datos, en la que se almacenan, entre otros, los siguientes elementos particulares:
\begin{itemize}
    \item Microcontrolador principal.
    \item Marca y modelo del módem LoRaWAN.
    \item Tipo de antena.
\end{itemize}

\subsubsection{Circuito de aplicación}

Cada posible configuración eléctrica tendrá asignado un identificador de catálogo (CatID), repetible solamente entre el mismo tipo de circuito. De tal manera, cada cambio que altere la funcionalidad del \emph{hardware}, modificando el circuito, requerirá obligatoriamente un nuevo CatID. Entre los cambios notables, se presentan los siguientes:
\begin{itemize}
    \item Agregar/retirar componentes.
    \item Cambiar valores de componentes alterando las impedancias.
    \item Cambiar las posiciones relativas de los conectores.
\end{itemize}

Cambios en el encaminamiento de pistas o en la orientación de componentes discretos que no afecten el funcionamiento de los dispositivos generalmente no implicarán la creación de un nuevo CatID.

%%%%%%%%%%%%%%%%%%%%%%%%%%%%%%%%%%%%%%%%%%%
\section{Abstracción de interfaces eléctricas}

En el mercado existe gran variedad de sensores, cuyas variables de medida pueden reflejarse análogamente en propiedades eléctricas como algún voltaje, corriente, impedancia, o alguna combinación en particular. Del mismo modo, estos sensores, que incluyen gran parte de la cadena de medida y que presentan la variable mediante una interfaz digital, tienen su propio orden para la carga útil de cada una de estas interfaces.

Se pretende que el procesamiento de la información en el nodo sea mínimo, con el fin de aprovechar al máximo la energía disponible, sobre todo al venir de baterías o fuentes no disponibles la mayor parte del tiempo, trasladando el procesamiento intensivo a quienes hagan uso de la aplicación a la cual se enviarán los datos de los sensores, casi en crudo.

\chapter{Sensores}
La medición de la calidad del aire se apoya en la evaluación y cuantificación de la presencia y concentración de contaminantes y componentes atmosféricos que pueden afectar la salud humana, el medio ambiente y los ecosistemas. Esta medición se realiza utilizando una variedad de sensores especializados para monitorear diferentes parámetros y contaminantes presentes en el aire. La medición de la calidad del aire se realiza a través de redes de monitoreo, que consisten en estaciones fijas distribuidas estratégicamente en diferentes ubicaciones geográficas. También se utilizan sensores portátiles y dispositivos de monitoreo en tiempo real para obtener mediciones más detalladas y específicas en áreas particulares o durante estudios de campo.

El análisis de los datos recolectados permite evaluar la calidad del aire, identificar fuentes de contaminación, evaluar la eficacia de las medidas de control y tomar decisiones informadas para proteger la salud y el medio ambiente. Además, estos datos son fundamentales para establecer políticas y regulaciones ambientales que promuevan una mejor calidad del aire y un entorno más saludable.

\section{Clasificación de sensores para calidad del aire}

\textls[-10]{Entre las principales variables que se evalúan dentro de lo que se considera calidad del aire, se encuentran los contaminantes gaseosos, las partículas suspendidas, los contaminantes biológicos y los parámetros meteorológicos.}

\subsection{Contaminantes gaseosos}

Dentro de esta denominación se incluyen gases como el dióxido de carbono (CO\textsubscript{2}), óxidos de nitrógeno (NO\textsubscript{2}), el ozono (O\textsubscript{3}), el dióxido de azufre (SO\textsubscript{2}) y algunos compuestos orgánicos volátiles (COV), entre otros. Estos contaminantes se miden para determinar su concentración y evaluar si están dentro de los límites aceptables establecidos por las normativas ambientales. Un esquema con los elementos y los principios básicos de su medición se muestra en la figura \ref{fig:fig1A}.

\begin{figure}[ht]
  \begin{center}
    \caption{Esquema básico de un sensor para contaminantes gaseosos \cite{sengascon}}
    \label{fig:fig1A}
    \includegraphics[width=.8\textwidth]{Fig1A}
    \propia
  \end{center}
\end{figure}

Los principales sensores para contaminantes gaseosos incluyen:

\begin{itemize}
  \item Sensor de dióxido de carbono (\CO): mide la concentración de dióxido de carbono en el aire. Es utilizado para controlar la ventilación en edificios, monitorear la calidad del aire en interiores y detectar posibles fuentes de contaminación.
  \item Sensor de COV: mide la presencia de compuestos químicos volátiles en el aire, como los emitidos por productos químicos, solventes, productos de limpieza y materiales de construcción. Los COV pueden contribuir a la mala calidad del aire y afectar la salud humana.
  \item Sensor de óxidos de nitrógeno (\NOX): mide la concentración de óxidos de nitrógeno en el aire, incluyendo el dióxido de nitrógeno ($\textrm{NO}_2$) y el monóxido de nitrógeno ($\textrm{NO}$). Estos gases son emitidos principalmente por la combustión de combustibles fósiles y pueden tener efectos negativos en la calidad del aire y la salud.
  \item Sensor de ozono ($\textrm{O}_3$): mide la concentración de ozono en el aire. El ozono puede ser perjudicial para la salud humana en altas concentraciones y es un contaminante común en áreas urbanas y cerca de fuentes de emisión.
\end{itemize}

\subsection{Partículas suspendidas}

Corresponde a la concentración de partículas, tanto finas (PM2.5) como grandes (PM10), en el aire. Estas partículas pueden provenir de fuentes naturales o antropogénicas, como la quema de combustibles fósiles, incendios forestales, el polvo de construcción, entre otros. Dentro de la familia de sensores para la medición de este tipo de afectación del aire, se tienen los siguientes:

\begin{itemize}
  \item Sensor de dispersión láser: utiliza un haz de luz láser para medir el tamaño y la concentración de partículas suspendidas. Mide la dispersión de la luz causada por las partículas en el aire y proporciona información sobre el tamaño de estas y su distribución.
  \item Sensor de impacto: este tipo de sensor utiliza una placa o sustrato sensible al impacto de las partículas. Cuando las partículas chocan contra la superficie del sensor, se genera una señal eléctrica proporcional a la concentración de partículas.
  \item Sensor de efecto túnel: este tipo de sensor se basa en el principio del efecto túnel cuántico, en el que las partículas suspendidas en el aire pueden alterar la corriente que fluye a través de una barrera de túnel. La alteración de la corriente se detecta y se utiliza para medir la concentración de estas.
  \item Sensor de masa: este sensor utiliza un filtro o un medio poroso para atrapar las partículas suspendidas en el aire. Luego, la masa de las partículas capturadas se pesa y se determina la concentración de estas en función de la masa medida.
  \item Sensor de condensación de vapor: este sensor utiliza la técnica de condensación de vapor para aumentar el tamaño de las partículas y hacerlas detectables. El vapor se condensa sobre las partículas, haciéndolas crecer en tamaño, y luego se mide el tamaño de las partículas condensadas.
\end{itemize}

\subsection{Contaminantes biológicos}

Se pueden medir microorganismos como bacterias, hongos, polen y esporas, entre otros microorganismos, que pueden representar riesgos para la salud humana. Algunos de los tipos de sensores utilizados para la detección de contaminantes biológicos son:

\begin{itemize}
  \item Sensores de PCR (reacción en cadena de la polimerasa): utilizan la técnica de PCR para amplificar y detectar material genético específico de los organismos objetivo. Estos sensores son altamente sensibles y se utilizan comúnmente para la detección de virus, bacterias y otros patógenos.
  \item \textls[-10]{Sensores de inmunodetección: utilizan anticuerpos u otras mo\-lécu\-las de reconocimiento específicas para detectar la presencia de agentes biológicos. Estos sensores pueden detectar proteínas o componentes celulares específicos y se utilizan en aplicaciones como la detección de toxinas bacterianas o virus.}
  \item Sensores de microarrays de ADN: utilizan una matriz de sondas de ADN específicas para detectar y distinguir diferentes secuencias de ADN. Estos sensores son útiles para la identificación de microorganismos y pueden detectar múltiples especies o cepas en una sola prueba.
  \item Sensores de espectroscopia: utilizan técnicas de espectroscopia para detectar y analizar la firma espectral de los contaminantes biológicos. Estos sensores pueden utilizar diferentes rangos del espectro electromagnético, como la luz visible, infrarroja o ultravioleta, para identificar características específicas de los organismos o sus productos de desecho.
  \item Sensores basados en biosensores: utilizan sistemas biológicos, como enzimas o células vivas, para detectar la presencia de contaminantes biológicos. Estos sensores pueden detectar cambios en la actividad enzimática o respuestas biológicas específicas para identificar la presencia de microorganismos o sustancias biológicas.
\end{itemize}

\subsection{Parámetros meteorológicos}

Además de los contaminantes específicos, también se miden variables meteorológicas como la temperatura, la humedad relativa y la velocidad y dirección del viento. Estos datos son importantes para comprender los patrones de dispersión de los contaminantes y su impacto en diferentes áreas. Algunos sensores destacados para estas magnitudes son:

\begin{itemize}
  \item Sensor de temperatura: mide la temperatura ambiente y puede utilizar termistores, termopares, sensores de resistencia (RTD), sensores semiconductores y sensores metálicos.
  \item Sensor de humedad relativa: mide la cantidad de humedad en el aire, de forma indirecta, por medición directa de la capacitancia, la resistencia o la conductancia eléctrica del medio húmedo.
  \item Sensor de presión atmosférica: mide la presión atmosférica, indirectamente, utilizando medios actinométricos sobre cuerpos de prueba como membranas, fuelles, tubos, entre otros.
  \item Sensores de lluvia o pluviómetros: miden la cantidad de precipitación que cae en un determinado periodo de tiempo. Los pluviómetros generalmente consisten en un recipiente graduado que recoge y mide el agua de lluvia, en un área y un tiempo determinados, para dar base a la proyección de cálculos asociados al volumen y caudal totales obtenibles de dichos parámetros.
\end{itemize}

\section[Sensores específicos para contaminantes gaseosos]{Sensores específicos para\\ contaminantes gaseosos}

\subsection{Sensores para CO\textsubscript{2}}

Las técnicas de detección de \CO\ han recorrido un largo camino en el desarrollo de sensores para el análisis de gases, pasando por métodos químicos y físicos como la cromatografía, el análisis infrarrojo, la espectrometría, la conductimetría, entre otras \cite{Kumar-2007}.

Actualmente se manejan, fundamentalmente, dos tipos de técnicas para la detección de \CO:
\begin{itemize}
  \item Determinación de la concentración de CO\textsubscript{2} en el aire por el método \gls{ndir}, el cual es el más empleado para la detección en tiempo real.
  \item Detección por medio de sensores electrolíticos de estado sólido.
\end{itemize}

\subsubsection{Sensor infrarrojo no dispersivo}
Un sensor \gls{ndir} es un dispositivo espectroscópico simple de uso frecuente como detector de gas.
Los componentes principales de este sensor se presentan en la figura \ref{fig:fig2A} y son los siguientes \cite{Yi-2005}:
\begin{itemize}
  \item Una lámpara de IR que se constituye en la fuente \gls{ndir}. Es el sistema \gls{ndir} propiamente dicho.
  \item Una cámara o tubo de luz por el cual circula la muestra.
  \item Un filtro de IR.
  \item Un detector de luz IR.
\end{itemize}

El gas se difunde en la cámara y la concentración de \CO\ se mide ópticamente por la absorción de una longitud de onda específica en el rango del infrarrojo. La luz IR se dirige, atravesando la cámara que contiene la muestra, hacia el detector; un filtro óptico, en frente del detector, elimina todas las componentes (longitudes de onda) de la luz, excepto la longitud de onda que absorben las moléculas del gas seleccionado. El \CO\ presenta una fuerte absorbancia en la banda IR; la longitud de onda frecuentemente empleada es de \SI{4260}{nm}.

\begin{figure}[H]
  \begin{center}
    \caption{Sensor \gls{ndir} para \CO}
    \label{fig:fig2A}
    \includegraphics[width=0.9\textwidth]{Fig2A}
    \propia
  \end{center}
\end{figure}

El término no dispersivo se refiere a que toda la luz (todas las longitudes de onda) atraviesa la muestra de gas y el filtrado se aplica solamente a la luz que logra llegar al final de la cámara, que es cuando esta ingresa al detector.

\subsubsection{Función característica}

La absorción de la radiación IR por parte de las moléculas de \CO\ sigue la ley de Beer-Lambert, que se expresa así:
\begin{align}
  A&=\log_{10}(I_a/I_f)\\
  A&=a \times x \times C
  \label{Lamb}
\end{align}
Donde:

\begin{itemize}
  \item[$A$] $\rightleftharpoons$ absorbancia del \CO. Es una magnitud adimensional.
  \item[$I_a$] $\rightleftharpoons$ Intensidad de la luz resultante en cualquier punto del sistema a lo largo de la cámara. Se mide en W/sr.
  \item[$I_f$] $\rightleftharpoons$ Intensidad de la luz que emerge de la fuente. Se mide en W/sr.
  \item[$a$] $\rightleftharpoons$ Coeficiente de absorción del \CO; se denomina igualmente como \emph{absortividad molar}. Se mide en L/(mol$\cdot$cm).
  \item[$x$] $\rightleftharpoons$ Longitud de la trayectoria, camino óptico recorrido o longitud de la celda. Se mide en cm.
  \item[$C$] $\rightleftharpoons$ Concentración del \CO. Se mide en mol/L.
\end{itemize}

La intensidad de luz, $I$, que llega al sensor óptico, disminuye en la medida en que se aumenta la concentración, $C$, de \CO\ en la cámara. Si la longitud de paso óptico de la cámara es $L$ y la absortividad molar del \CO\ es $A$, la intensidad de luz en el sensor óptico queda definida por la ley de Beer-Lambert:
\begin{equation*}
  I=I_0\times 10^{CLA}
\end{equation*}
Donde $I_0$ es la luz a la salida de la fuente.

\subsection{Sensores para COV}
Un sensor de COV es ampliamente utilizado en aplicaciones de monitoreo ambiental, control de calidad del aire y detección de gases en cualquier entorno.
El funcionamiento de un sensor de COV se basa en la interacción de los compuestos orgánicos volátiles con una capa sensible en el sensor. Cuando los COV están presentes en el aire, interactúan con la capa sensible, lo que produce cambios en sus propiedades eléctricas, ópticas o químicas. Estos cambios son medidos por el sensor y se convierten en una señal eléctrica que puede ser interpretada como la concentración de COV presente en el ambiente.

El modelo matemático de un sensor de COV varía según el tipo de sensor y la tecnología empleada. Algunos sensores de COV utilizan resistencias eléctricas, conductancias iónicas o espectroscopia para medir los cambios en la denominada capa sensible, y así poder calcular la concentración de COV.
Normalmente, cada fabricante de sensores para COV puede tener su propio modelo matemático, basado en el diseño y la tecnología utilizados en su sensor específico. Los modelos más detallados pueden involucrar ecuaciones racionales o empíricas, basadas en parámetros o en la respuesta del sistema, respectivamente.
El sensor de COV consta de varios componentes principales que trabajan juntos para detectar y medir la concentración de COV en el aire. Estos componentes son los siguientes:
\begin{itemize}
  \item La capa sensible: es la parte del sensor que entra en contacto con los COV presentes; está diseñada para interactuar con los COV y experimentar cambios en sus propiedades físicas o químicas cuando se expone a ellos.
  \item El transductor: es el elemento que convierte los cambios en las propiedades de la capa sensible en una señal eléctrica medible. Dependiendo del tipo de sensor, el transductor puede ser un resistor, un condensador, un diodo o algún otro dispositivo que responda a los cambios en la capa sensible.
  \item Un circuito de acondicionamiento de señal: es el circuito electrónico que amplifica y procesa la señal eléctrica generada por el transductor. Este circuito puede incluir amplificadores, filtros y otros componentes para garantizar una lectura precisa y estable de la concentración de COV.
  \item La fuente de alimentación: proporciona la energía necesaria para el funcionamiento de todo el sistema.
  \item La interfaz de salida: es la parte del sensor que permite la conexión con otros dispositivos o sistemas para transmitir la información sobre la concentración de COV. Esto puede incluir interfaces analógicas, como voltaje o corriente, o interfaces digitales, como un bus de comunicación serial.
\end{itemize}
Es importante destacar que los componentes específicos y la configuración del sensor pueden variar según el diseño y la tecnología utilizada. Cada fabricante puede tener su propia implementación y selección de componentes para lograr el mejor rendimiento y precisión en la medición de COV.

\subsection{Sensores electroquímicos para \texorpdfstring{NO\boldmath{$_{x}$}}{NOx}}

Los sensores electroquímicos para óxido de nitrógeno (\NOX) se basan en reacciones electroquímicas que ocurren en una celda sensora. Estos sensores constan de tres componentes principales: un electrodo de trabajo, un electrodo de referencia y un contraelectrodo. A continuación se presenta un modelo matemático híbrido (racional-empírico) utilizado en este tipo de sensores.

La ley de Nernst establece la relación entre el potencial eléctrico y la concentración de especies químicas en una celda electroquímica. El modelo matemático para un sensor de \NOX\ puede ser expresado así:
%%%%inicioecuacion
\begin{gather}
  \label{nox}
  E=E_o-(RT/nF)\times \ln(C)
\end{gather}
Donde:

\begin{itemize}
  \item[$E$] $\rightleftharpoons$  Potencial eléctrico medido en el electrodo.
  \item[$E_o$] $\rightleftharpoons$  Potencial de base de la reacción electroquímica.
  \item[$R$] $\rightleftharpoons$  Constante universal de los gases ideales.
  \item[$T$] $\rightleftharpoons$  Temperatura termodinámica.
  \item[$F$] $\rightleftharpoons$  Constante de Faraday.
  \item[$C$] $\rightleftharpoons$  Concentración del \NOX.
\end{itemize}

Este modelo establece una relación logarítmica entre el potencial eléctrico medido y la concentración de \NOX.
Existe, también, un modelo denominado de difusión y reacción; en dicho modelo se considera la difusión de los gases de \NOX\ desde el entorno hacia la superficie del electrodo, donde ocurren las reacciones electroquímicas. El modelo matemático puede incluir las ecuaciones de difusión de Fick y reacciones de oxidación y reducción en el electrodo para describir el comportamiento del sensor.
Otros modelos más complejos pueden considerar parámetros adicionales, como el coeficiente de difusión, la velocidad de transferencia de electrones y la resistencia interna del sensor, para obtener una descripción más precisa del comportamiento del sensor electroquímico para \NOX.
Los modelos matemáticos específicos pueden variar según el diseño y la tecnología del sensor electroquímico utilizado. Los fabricantes de sensores electroquímicos suelen proporcionar documentación técnica que incluye información sobre el modelo matemático y las ecuaciones asociadas utilizadas en sus sensores específicos.

\section{Sensores específicos para material particulado}

\subsection{Sensor de dispersión láser para material particulado}

El sensor de dispersión láser para material particulado mantiene la técnica de dispersión de un haz de luz que atraviesa el material por analizar, sufriendo diferentes niveles de atenuación en sus componentes de frecuencia, debido a las diferentes partículas encontradas en su camino. Los elementos que lo componen y sus funciones individuales se pueden resumir así:
\begin{itemize}
  \item Fuente láser: puede ser un láser, generalmente de infrarrojo, con un espectro amplio.
  \item Cámara de dispersión: permite el encuentro entre el láser y el aire contaminado con las partículas en cuestión.
  \item Detector de partículas: los métodos comunes incluyen la dispersión de luz debido al encuentro entre partículas y el haz de luz.
  \item Contador de partículas: una vez que las partículas son detectadas, el sensor captura y cuenta las partículas presentes en una muestra de aire. Esto se puede lograr mediante técnicas como la atracción electrostática, la filtración o la interacción de las partículas con una superficie sensible.
  \item Conversor de partículas a señal eléctrica: el sensor convierte la información de las partículas detectadas en una señal eléctrica, generalmente corriente.
\end{itemize}
Un sensor típico de esta tecnología es el producido por Honeywell, cuyo esquema simplificado se presenta en la figura \ref{fig:fig3A}.
Es importante tener en cuenta que los sensores de material particulado pueden utilizar diferentes tecnologías y principios de funcionamiento, como la dispersión láser, la captación en filtros o la medición de la conductividad.

Un posible modelo matemático para un sensor de partículas de dispersión de luz es el siguiente:
\begin{align}
  C = k \times (I / D)
\end{align}

Donde:
\begin{itemize}
  \item[$C$] $\rightleftharpoons$ Concentración de partículas en el aire.
  \item[$I$] $\rightleftharpoons$ Intensidad de la luz dispersada o atenuada por las partículas.
  \item[$D$] $\rightleftharpoons$ Distancia entre la fuente de luz y el detector.
  \item[$k$] $\rightleftharpoons$ Sensibilidad de la cámara de dispersión.
\end{itemize}

Este modelo asume una relación lineal entre la intensidad de la luz dispersada y la concentración de partículas. La sensibilidad de la cámara se determina experimentalmente y puede variar dependiendo de las características específicas del sensor y del tipo de partículas que se estén midiendo.

\begin{figure}[H]
  \centering
  \caption{Sensor de dispersión para material particulado}
  \label{fig:fig3A}
  \includegraphics[width=0.7\textwidth]{Fig3A}

  \propia
\end{figure}

\subsubsection[Sensor de dispersión láser multiángulo para material particulado]{Sensor de dispersión láser multiángulo\\ para material particulado}

Un novedoso sensor de dispersión para partículas PM2.5 es el desarrollado en el State Key Laboratory of Industrial Control Technology, del College of Control Science and Engineering, de la Universidad de Zhejiang, el cual permite eliminar los efectos de factores meteorológicos, como la humedad, en las partículas. Se concentra en detectar la dispersión en tres ángulos diferentes: \SI{40}{\degree}, \SI{55}{\degree} y \SI{140}{\degree}. Un esquema con arquitectura similar a la disposición real del sensor se presenta en la figura \ref{fig:fig4A}.

\begin{figure}[h]
  \centering
  \caption{Sensor de dispersión multiángulo para material particulado}
  \label{fig:fig4A}
  \includegraphics[width=\textwidth]{Fig4A}

  \propia
\end{figure}

\section{Interfaces de sensores}

\subsection{Sensores con interfaz análogica}

Cuando un sensor presenta un cambio en sus propiedades eléctricas proporcional a los cambios en su variable de medida, es recomendable implementar una cadena de medida que permita su correcto funcionamiento y a su vez un aprovechamiento máximo del rango completo de detección del \gls{adc} del microcontrolador. También deberá considerarse si es necesario hacer algún tipo de tratamiento básico a los datos (promedio, filtrado, entre otros) antes de concluir un valor que se pueda enviar por la red LoRaWAN.

\subsection{Sensores con interfaz digital}

La estructura general consiste en el envío de uno o más comandos, con el fin de obtener una respuesta a estos. El formato de cada comando y su respuesta asociada depende fuertemente del fabricante de cada sensor en particular.

\chapter{Concentrador}

\section{Hardware}
El hardware asociado a el concentrador contiene dos elementos: 

\begin{itemize}
    \item Dispositivo anfitrión (Host)
    \item Núcleo del concentrador
\end{itemize}

\subsection{Host}
El host contiene el dispositivo computacional, que le permite al concentrador LoRaWAN controlar el dispositivo físico de comunicaciones. Es un computador embebido normalmente con sistema operativo, también embebido, tipo Linux, aunque podrán existir host más livianos, o básicos, que tengan un \gls{rtos} sobre un microprocesador ARM. Este dispositivo contiene software con un nivel de abstracción medio y alto donde enriquece las características de toda la red.
Las funcionalidades básicas del host son:

\begin{itemize}
    \item Redirección de paquetes LoRaWan una vez recibidos por el Núcleo del concentrador.
    \item Envío y recepción de información hacia y desde Internet a través de comunicaciones como Ethernet, Wi-Fi, telefonía móvil (2G/3G/4G/5G), enlace microondas o satelital, etc.
    \item Cifrado de la información para enviarla en redes no seguras como Internet.
    \item Comunicación orientada o no orientada a la conexión hacia el servidor LoRaWAN.
    \item Caché o estrategias de represamiento temporal de información, dadas las condiciones de red.
    \item Conexión reciliente frente a desconexiones del concentrador.
    \item cambio de parámetros de la red hacia el núcleo del concentrador, según los comandos enviados por el servidor LoRaWAN.
\end{itemize}

Cabe resaltar que los paquetes LoRaWAN ya tienen un cifrado punto a punto entre el dispositivo y el servidor LoRaWAN, el cifrado adicional que impone el host es entre este aparato y la red y/o servidor LoRaWAN, debido a que aparte de los paquetes LoRaWAN, también se transmite información propia del concentrador y no directamente relacionada con la comunicación escencial de la red de sensores o dispositivos clientes LoRaWAN.

Se entiende, entonces, que el host controla como un maestro al núcleo del concentrador y sirve como puente entre la modulación LoRa y la comunicación tradicional, conectando este concentrador a la red de redes (Internet) y junto con el papel del servidor LoRaWAN, generan una red orientada a INTERNET de las Cosas.

El host es una Raspberry Pi configurada de tal manera que su interfaz de pines \gls{gpio} se pueda conectar a el núcleo del concentrador.

\subsection{Núcleo del concentrador}
Es una tarjeta integrada, con referencia iC880A, como un transceiver diseñado para recibir simultáneamente paquetes LoRa a través de múltiples canales con diferentes \gls{sf}. Esta pasarela esta en la capacidad de comunicarse con miles de nodos a distancias lejanas, y al ser adaptado a un host, puede comunicarse con un servidor LoRaWAN para la acumulación y tratamiento de los datos. Esta pasarela esta en la banda de frecuencias de 868~MHz, con una sensibilidad de -138~dBm. Utiliza un modulo de radio del fabricante Semtech, el cual soporta varios anchos de banda y un rango de \gls{sf} de 7 a 12. La modulación LoRa de espectro ensanchado se realiza representado cada de información por múltiples chips de información. La velocidad con la que es enviada la información es referida a la \gls{rs} nominal. La relación entre la \gls{rs} y la rata de chip, representa el número de símbolos de modulación enviados por bit de información. El \gls{sf} se debe conocer de antemano por el nodo y la pasarela, ya que los diferentes \gls{sf} son ortogonales entre sí. 

\begin{figure}
\centering
\caption{Stack de Nodo, concentrador y red}
\includegraphics[width=\textwidth]{capitulos/figuras/lorawan-gateway-network-detailed.jpg}
\fuente{\url{https://www.i-scoop.eu/internet-of-things-guide/iot-network-lora-lorawan/}}
\label{f:stack_lora_detailed}
\end{figure}

\section{Software}

\subsection{Alpine Línux}
Alpine Linux es un sistema operativo muy liviano, para equipos computacionales muy sencillos y de uso general, está diseñada para usuarios avanzados que aprecian la seguridad, la simplicidad y la eficiencia de los recursos. Alpine línux se basa en uClibc y BusyBox. uClibc es una biblioteca en C diseñada para un sistema Linux embebido, con requerimientos de recursos de computo mínimos para arquitecturas x86, AMD64, ARM, Blackfin, H8300, m68k, MIPS, PowerPC, SuperH, SPARC y V850. BusyBox es un software que utiliza las ventajas de estándares de Unix en una sola consola.

\begin{quotation}
    Un contenedor no requiere más de 8 MB y una instalación mínima en el disco requiere alrededor de 130 MB de almacenamiento. No solo obtiene un entorno Linux completo, sino una gran selección de paquetes del repositorio.
    
    \fuente{\url{https://www.linuxadictos.com/disponible-la-nueva-version-de-alpine-linux-3-8-1.html\#google\_vignette}}
\end{quotation}

Los paquetes binarios se reducen y se dividen, lo que le brinda aún más control sobre lo que se instala, lo que a su vez mantiene su entorno lo más pequeño y eficiente posible.

El administrador de paquetes llamado apk, propio de Alpine línux, y el sistema de inicialización OpenRC, son configuraciones controladas por scripts. Esto le proporciona un entorno Linux simple y nítido sin todo el ruido, lo cual permite agregar solo los paquetes que se necesitan para el proyecto, ya sea que se esté construyendo un PVR doméstico o un controlador de almacenamiento iSCSI.

Alpine Linux fue diseñado pensando en la seguridad. Todos los binarios del área de usuario se compilan como ejecutables independientes de la posición (PIE) con protección contra la rotura de pilas. Estas funciones de seguridad proactiva evitan la explotación de clases enteras de vulnerabilidades de día cero y otras.

\subsubsection{Configuración de Alpine Linux}

Luego de seguir las indicaciones de la instalación en Raspberry PI , teniendo en cuenta que se debe dejar la partición FAT32 de al menos 500MB, también es recomendable usar el resto de espacio en la tarjeta SD como una partición ext4 para el espacio de cache, repos, etc.
Una vez instalado el sistema Alpine Linux, se deben hacer los siguientes procesos:

\begin{itemize}
    \item Configurar / colocar los archivos auxiliares de configuración externa en la partición FAT32 (/ media / mmcblk0p1).
    \item Crear un archivo .apkovl.tar.gz personalizado basado en la plantilla y los scripts asociados
    \item Subir el archivo .apkovl.tar.gz a la segunda partición (ext4) cuyo nombre sería / media / mmcblk0p2 para Alpine Linux.
    \item Configurar el proxy (si lo hay) con el comando \texttt{\$ setup-proxy} y luego importarlo en las variables de entorno locales \texttt{\$ source /etc/profile.d/proxy.sh}
    \item Ejecutar \texttt{\$ setup-apkcache} para crear el cache local de paquetes de Alpine Linux teniendo en cuenta que, en general, el cache quedará en \texttt{/media/mmcblk0p2/cache}
    \item Sincronizar el caché local con los repositorios mediante el comando \texttt{\$ apk cache -v sync}
\end{itemize}
    
Los archivos .apkovl.tar.gz son archivos de backup local de configuración que se usan para hacer más reproducible la configuración de los gateways.

\subsection{LoRa packet forwarder}

Es el software que se coloca en el host del gateway LoRa, el cual reenvía los paquetes de RF que recibe del concentrador a un servidor a través de un enlace IP/UDP.

Esta configuración está pensada para ejecutar imágenes de Alpine Linux, para usarlos con gateways IC880A en computadores/servidores como la Raspberry Pi (con un puerto USB 2.0, si la versión del IC880A es USB, o con puertos GPIO para las versiones IC880A SPI).

Se configura en el Gateway un software para redireccionamiento de paquetes LoRa llamado packet forwarder el cual utiliza un protocolo UDP (no orientado a la conexión y no cifrado) junto con un puente MQTT llamado LoRa Gateway Bridge el cual permite usar TCP+SSL para conectarse con la red LoRaWAN en vez de UDP. También puede emitir una señal Beacon sincrónica GPS en toda la red utilizada para coordinar todos los nodos de la red. 

\section{Software de configuración y despliegue}

EL modelo de red del proyecto, utiliza software libre para la conversión de protocolos (backhoul), almacenamiento de la información (Bases de datos) y de aplicación (información de usuario). Para estos procesos el servidor utiliza varias herramientas de software. Para la configuración de la aplicación se requiere de Lora Server, Influx DB y de grafana, las cuales se explicarán a continuación.  

\subsection{Lora server}

El servidor de red LoRaWAN es el núcleo del sistema de gestión de redes LoRa y constituye el componente central del proyecto LoRa Server. Su papel es fundamental para garantizar la operatividad y eficiencia de la red. Una de sus principales responsabilidades es la desduplicación de tramas de enlace ascendente, lo que significa que el servidor se encarga de procesar y filtrar las tramas de datos que llegan desde las puertas de enlace o gateways, evitando que se dupliquen y asegurando que solo se procese una versión de la información recibida.

\subsection{Influx DB}

InfluxDB es una base de datos especializada en series temporales, diseñada para gestionar grandes volúmenes de escritura y consultas de manera eficiente. Forma parte esencial de la pila TICK, que es un conjunto de herramientas para la recopilación, almacenamiento, visualización y alerta de datos en tiempo real. InfluxDB es ideal como repositorio para casos de uso que requieran almacenar y procesar grandes cantidades de datos con marcas de tiempo, como la monitorización en entornos DevOps, métricas de rendimiento de aplicaciones, datos provenientes de sensores de IoT y análisis en tiempo real.

Para integrar InfluxDB con LoRa Server, la conexión se realiza a través de LoRa App Server, donde se especifica la dirección del servidor donde está alojada la base de datos, junto con las credenciales de acceso (usuario y contraseña). Mediante scripts de JavaScript, el servidor es capaz de crear automáticamente las tablas y series temporales necesarias para almacenar los datos que se ingresen en la base de datos. Esta capacidad automatizada facilita la gestión de grandes volúmenes de datos y permite una interacción fluida entre los dispositivos IoT y la base de datos para una supervisión y análisis en tiempo real.

\subsection{Grafana}

Grafana es una herramienta de código abierto, licenciada bajo Apache 2.0, diseñada para la visualización y el formateo de datos métricos. Ofrece la capacidad de crear paneles interactivos y gráficos informativos a partir de diversas fuentes de datos, incluyendo bases de datos de series temporales como Graphite, InfluxDB y OpenTSDB.

En primera instancia se instala el software en el servidor, para luego configurarlo con la base de datos utilizada. Para la visualización de los datos se configuran las consultas a la base de datos en cada uno de los Dashboards disponibles con la forma de visualización deseada.

\chapter[Aplicación propuesta]{Aplicación propuesta: medición de la calidad del aire} \label{cap:aplicacion}

La calidad del aire afecta la salud de los seres vivos y el medio ambiente. Entre los efectos asociados con la polución del aire, se hallan la disminución de la función pulmonar, las enfermedades cardiovasculares, el incremento de incidentes asmáticos, entre otros. Además, del impacto nocivo en el ser humano, se afecta la vegetación, decrece la visibilidad y hay una alta influencia en el clima global.

\section{Contaminación del aire}

La contaminación del aire es la presencia de una mezcla compleja de diferentes compuestos químicos en forma de partículas sólidas, gotas líquidas o gases. Algunos de estos contaminantes son de corta duración en la atmósfera, mientras que otros son de larga vida. La cantidad de tiempo que un contaminante particular permanece en la atmósfera depende de su afinidad reactiva con otras sustancias y/o su tendencia a depositarse en una superficie y de las condiciones climáticas, incluidas la temperatura, la luz solar, la precipitación y la velocidad del viento.

Los contaminantes son emitidos por una amplia variedad de fuentes naturales y artificiales. Algunos ejemplos de fuentes artificiales son las plantas generadoras de electricidad, los automóviles y las plantas de producción de petróleo y gas. Las fuentes de contaminantes naturales incluyen incendios forestales, tormentas de polvo y erupciones volcánicas, entre otras. Los denominados \emph{contaminantes primarios} son emitidos directamente de la fuente y están representados por material particulado (PM2.5; PM10), el monóxido de carbono (CO), el dióxido de nitrógeno (NO\textsubscript{2}), el dióxido de azufre (SO\textsubscript{2}) y el plomo (Pb). Los conocidos como \emph{contaminantes secundarios} se obtienen de reacciones químicas; este grupo incluye el ozono (O\textsubscript{3}) y algunos tipos de material particulado. Las concentraciones de contaminantes en el aire varían en el espacio y en el tiempo debido a otros parámetros, a la proximidad con las fuentes y a las condiciones climáticas.

\subsection{Efectos en la salud}

La \gls{epa} ha establecido el \gls{aqi} para trasladar las medidas de polución en potenciales efectos en los individuos. El \gls{aqi} es un indicador que reporta la calidad diaria del aire e indica qué tan limpio o contaminado está y qué efectos implica para la salud.

La \gls{epa} calcula el \gls{aqi} basado en los cinco mayores contaminantes del aire: nivel de ozono, polución de partículas, monóxido de carbono, dióxido de azufre y dióxido de nitrógeno.

La tabla~\ref{tabla:epa} describe los efectos en la salud de acuerdo al \gls{aqi}.

\begin{table}[H]
  \centering
  \caption{Estándares Nacionales de Calidad del Aire Ambiental (NAAQS)}
  \label{tabla:epa}
  \footnotesize
  \begin{tblr}{
      width = \textwidth,
      colspec = {llX[m]l},
      hlines,
      vspan = even,
    }
    \textbf{Contaminante} & \textbf{Primario/secundario} & \textbf{Tiempo promedio} & \textbf{Nivel (límite)} \\

    \SetCell[r=2]{l} Monóxido de carbono (CO)
    & Primario & 8 horas & 9 ppm \\
    & Primario & 1 hora & 35 ppm \\

    Plomo (Pb)
    & Primario y secundario & Promedio trimestral móvil & 0.15 µg/m$^3$ \\

    \SetCell[r=2]{l} Dióxido de nitrógeno (NO$_2$)
    & Primario & 1 hora & 100 ppb \\
    & Primario y secundario & Anual & 53 ppb \\

    Ozono (O$_3$)
    & Primario y secundario & 8 horas & 0.070 ppm \\

    \SetCell[r=2]{l} Partículas (PM2.5)
    & Primario & Anual & 9.0 µg/m$^3$ \\
    & Primario y secundario & 24 horas & 35 µg/m$^3$ \\

    Partículas (PM10)
    & Primario y secundario & 24 horas & 150 µg/m$^3$ \\

    \SetCell[r=2]{l} Dióxido de azufre (SO$_2$)
    & Primario & 1 hora & 75 ppb \\
    & Secundario & 3 horas & 0.5 ppm \\
  \end{tblr}

  \vspace*{.3\baselineskip}

  \fuente{EPA}
\end{table}

\section{Monitoreo de contaminantes del aire}

La \gls{epa} ha identificado seis contaminantes estándar de alto nivel de importancia para la salud y el medio ambiente: ozono (O\textsubscript{3}), material particulado (PM), monóxido de carbono (CO), dióxido de nitrógeno (NO\textsubscript{2}), dióxido de azufre (SO\textsubscript{2}) y plomo (Pb). La \gls{epa} ha establecido los \gls{naaqs}, los cuales clasifica en primarios y secundarios; los primarios están destinados a proteger la salud pública, y los secundarios, a proteger el bienestar público.

Con la nueva generación de sensores de calidad del aire, de bajo costo y portátiles, se tienen muchas facilidades para su uso en aplicaciones de esta área. Aunque los sensores de contaminación del aire aún se encuentran en su etapa inicial de desarrollo, la \gls{epa} tiene pautas específicas que se deben aplicar para establecer un monitoreo de aire reglamentario. La tabla~\ref{tabla:epa1} resume algunas áreas potenciales de aplicación no regulatoria para sensores de aire y proporciona descripciones breves y ejemplos.

\begin{table}[H]
  \centering
  \caption{Áreas de aplicación no regulatoria para sensores de aire}
  \label{tabla:epa1}
  \footnotesize
  \begin{tblr}{
      width = \textwidth,
      colspec = {Q[l,3cm] X[l] X[l]},
      hlines,
      vspan = even,
    }
    \textbf{Área de aplicación} & \textbf{Descripción breve} & \textbf{Ejemplos de uso} \\

    Educación y ciencia participativa
    & Uso en entornos escolares o proyectos comunitarios para fomentar el aprendizaje sobre la calidad del aire.
    & Lecciones STEM en escuelas; proyectos de ciencia ciudadana para mapear aire local. \\

    Investigación científica
    & Recopilación de datos de alta resolución temporal o espacial para estudios académicos o de salud.
    & Estudios epidemiológicos; verificación de modelos de dispersión de aire. \\

    Información local y comunal
    & Entrega de datos específicos de una ubicación que no cuenta con estaciones de monitoreo oficiales.
    & Identificación de ``puntos calientes'' (\emph{hotspots}) de contaminación en barrios específicos. \\

    Respuesta a emergencias
    & Monitoreo rápido durante eventos inesperados que afectan la calidad del aire de forma temporal.
    & Medición de humo durante incendios forestales o accidentes industriales. \\

    Monitoreo de exposición personal
    & Evaluación del aire que una persona respira mientras se desplaza en su vida diaria.
    & Uso de sensores portátiles mientras se camina, se viaja en carro o se está en el trabajo. \\

    Identificación de fuentes de emisión
    & Ayuda en la localización de fugas o emisiones no controladas en instalaciones cercanas.
    & Detección de fugas en plantas industriales o monitoreo cerca de chimeneas sospechosas. \\

    Monitoreo de aire interior
    & Evaluación de contaminantes dentro de edificios o residencias.
    & Evaluación de la ventilación en oficinas o detección de PM2.5 por estufas de leña. \\
  \end{tblr}
  
  \vspace*{.3\baselineskip}

  \fuente{EPA}
\end{table}

\section{Modelo de red para la aplicación}

El modelo de red implementado obedece a la topología estrella. El funcionamiento de esta red está soportado y compuesto por nodos, pasarelas y un servidor de aplicaciones (figura \ref{f:modelo}).

\begin{figure}[H]
  \centering
  \caption{Modelo de la red LoRaWAN}
  \label{f:modelo}
  \includegraphics[width=\textwidth]{M_lora}
  \propia
\end{figure}

Cada nodo enviará a la red un conjunto de datos geoposicionados, compuesto por la medición de algunos parámetros del aire a través de sensores de temperatura, humedad, dióxido de carbono, presión atmosférica y material particulado.

La red está compuesta por tres \emph{hosts}-pasarelas ubicadas en las siguientes sedes de la Universidad Distrital: Ingeniería, Aduanilla de Paiba y la sede de posgrados. La red cuenta inicialmente con cinco nodos de prueba.

Cada pasarela de la red tiene ocho puertos de atención a los nodos y otros ocho puertos tipo \gls{rfu}. En el caso de la configuración en la clase A, una pasarela está en la capacidad de atender a más de 5000 nodos en su huella de acción.

\section{Cobertura}

Inicialmente, las pasarelas de la red \gls{lwan} implementada se colocaron en la sede de la Facultad de Ingeniería y en la sede de la Facultad Tecnológica, siendo dichas pasarelas uno de los más importantes objetos del proyecto. Un indicador fundamental del proyecto fue la cobertura, la cual se evaluó ubicando en cada nodo un elemento de geoposicionamiento global, donde la red \gls{lwan} toma los datos de cada nodo y los lleva a la nube. En el proyecto se modelaron tanto pasarelas como nodos con núcleos del fabricante ISMT. También se dispuso de algunos elementos de red comerciales (Multitech e ISMT) para comprobar la interoperabilidad.

Los datos fueron procesados utilizando herramientas de analítica para luego ser mostrados en la aplicación gráfica Grafana. Las figuras de esta sección muestran los resultados de las coberturas evaluadas en la red propuesta.

\subsection{Cobertura en la sede Tecnológica}

\subsubsection{Prueba de interoperabilidad}

Una de las pasarelas del fabricante Multitech, denominada Mote II, se evaluó con el nodo del fabricante ISMT. Lo primero que se quería comprobar era la interoperabilidad, para lo cual se instaló en la sede Tecnológica, y se modelaron los \emph{plugins} que reconocen el lenguaje de los \emph{logs} entre la pasarela y el nodo.

\begin{figure}[H]
  \centering
  \caption[Cobertura e interoperabilidad Multitech en inmediaciones del campus]{Cobertura e interoperabilidad Multitech\\ en inmediaciones del campus}
  \label{f:mapa1}
  \includegraphics[width=.6\textwidth]{cobertura/fig1.png}

  \propia
\end{figure}

En la figura \ref{f:mapa1} se observa la cobertura en el interior de la sede Tecnológica, con el \emph{gateway} Multitech y el nodo Mote II. El resultado muestra cobertura completa en la sede tecnológica, lo que significa que se podría implementar un campus inteligente.

De la misma forma, en la figura \ref{f:mapa2} se presenta la interoperabilidad de la pasarela modelada en este proyecto con un nodo del fabricante ISMT, denominado Mote II. Esto muestra que tiene prestaciones similares a las de los fabricantes certificados.

\begin{figure}[H]
  \centering
  \caption[Cobertura e interoperabilidad  Mote II en inmediaciones del campus]{Cobertura e interoperabilidad Mote II\\ en inmediaciones del campus}
  \label{f:mapa2}
  \includegraphics[width=.6\textwidth]{cobertura/fig2.png}

  \propia
\end{figure}

En la figura \ref{f:mapa2} se muestra la cobertura exterior, con la pasarela Multitech y el nodo del proyecto. El resultado es un alcance aproximado de 200 m de cobertura de la red instalada en esta sede.

El corto alcance se debe a que no se tiene una torre lo suficientemente alta para evitar obstáculos entre la pasarela y el nodo. Todas las diferentes configuraciones mostraron las mismas características en cobertura.

Estos valores de cobertura están alejados de las normas de la alianza \mbox{\gls{lwan}}, las cuales establecen, para zonas urbanas, algunos kilómetros de cobertura. También hay que tener en cuenta que la tecnología está aún en un periodo de evolución (2017, año de adquisición de los elementos tecnológicos), y que las certificaciones para fabricantes se iniciaron en el año 2016.

\subsubsection{Cobertura de la red modelada por el proyecto}

Una de las satisfacciones más importantes que deja el proyecto es su realización, dadas las limitaciones tecnológicas y la mesura en su presupuesto. La red \gls{lwan} del proyecto, como se mostró en la sección de modelamiento del nodo,  muestra una cobertura quizás con alcances levemente mayores que los de los fabricantes; esto se debe a que la antena del nodo fue realizada dentro del proyecto, lográndose una ganancia óptima.

\begin{figure}[H]
  \centering
  \caption{Cobertura de red modelada en el proyecto}
  \label{f:mapa3}
  \includegraphics[width=.6\textwidth]{cobertura/fig3.png}

  \propia
\end{figure}

En la figura \ref{f:mapa3} se muestra la cobertura exterior de la sede en mención, teniendo como resultado un alcance aproximado de 350\,m entre la pasarela y el nodo. La adquisición de datos fue hecha por estudiantes que hicieron recorridos por los sectores aledaños donde se instaló la pasarela.

\subsection{Cobertura en la sede de Ingeniería}

En la sede de Ingeniería se colocó una pasarela modelada con un chipset de radio LoRa del fabricante ISMT para el nodo y la plaqueta de radio para la pasarela. También se hicieron pruebas de interoperabilidad, cuyos resultados fueron muy satisfactorios dados los alcances obtenidos en este lugar en comparación con la red instalada en la sede Tecnológica. La mejora en el alcance se logró gracias a la pasarela, ya que se ubicó en la azotea del edificio administrativo, aproximadamente a cincuenta metros del suelo, lo cual permite que haya menos obstáculos y se cubra una mayor distancia.

\subsubsection{Red \gls{lwan} tipo malla}

En la sede de Ingeniería se implementó un modelo de red de malla, para lo cual se utilizaron las pasarelas del proyecto y la del fabricante Multitech, así como los nodos del proyecto y de los fabricantes Multitech e ISMT.

\begin{figure}[H]
  \centering
  \caption{Cobertura \emph{gateway} ISMT Mote II}
  \label{f:mapa4}
  \includegraphics[width=.7\textwidth]{cobertura/fig4.png}

  \propia
\end{figure}

En la figura \ref{f:mapa4} se muestra la cobertura exterior en la sede de Ingeniería; allí se aprecia la gran cantidad de datos sobre la posición de los nodos. En los datos se observó que efectivamente todas las pasarelas obtuvieron resultados, lo que significa que la red obedece a los protocolos de la alianza \gls{lwan}; es decir, si dos pasarelas reciben el mismo dato, se discrimina por el que haya llegado con mayor potencia. Una de las pasarelas fue ubicada en una azotea, y las otras dos en interiores (en la sala del grupo de investigación GITUD).

\subsubsection{Red \gls{lwan} en interiores}

Utilizando el modelo de red de malla, se realizó una discriminación para entornos interiores. Los resultados obtenidos en la sede de Ingeniería muestran que las pasarelas de interior fueron las que lograron captar datos, ya que la pasarela exterior se encontraba a una distancia considerable y con una gran cantidad de obstáculos.

\begin{figure}[H]
  \centering
  \caption{Cobertura en interiores}
  \label{f:mapa5}
  \includegraphics[width=.7\textwidth]{cobertura/fig5.png}

  \propia
\end{figure}

Para la lectura de datos en interiores, al realizar las mediciones, observó que puede penetrar obstáculos de hasta tres pisos, tal como se determinó en el edificio de Ingeniería. De acuerdo con los datos mostrados en la figura~\ref{f:mapa5}, se podría mejorar la cobertura en interiores aumentando la potencia o mejorando la ganancia de las antenas. Igualmente, se pueden adquirir elementos nuevos de radio LoRa que tengan mayor eficiencia y se acerquen a las condiciones dadas por los estándares de la alianza.

\section{Interpretación de los datos}

A continuación se presentan algunas figuras de las variables medidas, con las aclaraciones más relevantes y referentes a los objetivos del proyecto.

\subsection{Ciclos de enlace de los nodos}

Con los siete nodos disponibles, se efectuaron pruebas de comunicación continua para cada uno, las cuales duraron aproximadamente entre cinco y quince días por nodo. La figura \ref{gl1} muestra dichos ciclos de enlace durante un mes.

\begin{figure}[H]
  \centering
  \caption{Ciclos de trabajo por nodo}
  \label{gl1}
  \includegraphics[width=.7\textwidth]{Resultados/gl1.png}

  \propia
\end{figure}

Los nodos expuestos allí: IM880b-2, IM880b-l-01, IM880bl-03 y Pylates-5983, son productos de diseño y fabricación dentro del proyecto; los demás: Motell-0406, Motell-0405 y mDot-Box-01, son comerciales y adquiridos para el proyecto.

\subsection{Dióxido de carbono (\CO)}

Del nodo IM880-01 se presenta el comportamiento del medidor de \CO\ dentro del laboratorio GITUD del sexto piso del edificio Sabio Caldas. Los picos y las elevaciones de nivel que se suceden durante el día entre las seis de la mañana y las seis de la tarde corresponden a aquellos momentos en que el número de personas aumenta, con lo cual el nivel de \CO\ dentro de la sala se eleva por encima del nivel normal de referencia, cuando la sala está totalmente vacía (500 ppm). Las variaciones de la medida se presentan en la figura \ref{gl2} para un periodo de cuatro días continuos.

\begin{figure}[H]
  \begin{center}
    \caption{Medición de \CO\ dentro de un recinto}  \label{gl2}
    \includegraphics[width=.5\textwidth]{Resultados/gl2.png}

    \propia
  \end{center}
\end{figure}

\subsection{Iluminancia}

El nivel de iluminancia del laboratorio GITUD se aprecia en la figura \ref{gl3}. Allí se pueden diferenciar el encendido y apagado de las luminarias durante el día y las condiciones de 0\,lx durante la noche. La ubicación del sensor, en el nodo mDot-Box-01, con respecto a la luminaria hacía que este saturara en aproximadamente 4000\,lx.

\begin{figure}[H]
  \begin{center}
    \caption{Medición saturada de iluminancia}
    \label{gl3}
    \includegraphics[width=.7\textwidth]{Resultados/gl3.png}

    \propia
  \end{center}
\end{figure}

Para el sensor de iluminancia del nodo 82.light, la iluminación se atenuó y se graduó entre 50 lx y 70 lx, lo cual permite apreciar la precisión y resolución con que se realiza la medida. En la figura \ref{gl4} se muestran los resultados.

\begin{figure}[H]
  \begin{center}
    \caption{Medición de iluminancia dentro de un recinto}
    \label{gl4}
    \includegraphics[width=.45\textwidth]{Resultados/gl4.png}

    \propia
  \end{center}
\end{figure}

La figura \ref{gl6} muestra los cambios de iluminación en un recorrido a pie en exteriores. Los picos entre 10\,000 y 20\,000\,lx corresponden a exposiciones del medidor directamente enfrentando el sensor a la luz solar; los valles bajos son mediciones a la sombra de objetos o edificaciones. La medición se realizó con el IM880b-l-01, fabricado dentro del proyecto.

\begin{figure}[H]
  \begin{center}
    \caption{Medición de iluminancia en exteriores}\label{gl6}
    \includegraphics[width=.6\textwidth]{Resultados/gl6.png}

    \propia
  \end{center}
\end{figure}

\subsubsection{Latitud}

La figura \ref{gl5} corresponde a lecturas de confirmación de precisión de los datos de la variable latitud dentro de las instalaciones del GITUD, realizadas con el Motell/0405. Como se puede apreciar, las valoraciones mínimas (min), máxima (max) y promedio (avg) son iguales.

\begin{figure}[H]
  \begin{center}
    \caption{Medición de latitud}\label{gl5}
    \includegraphics[width =.6\textwidth]{Resultados/gl5.png}

    \propia
  \end{center}
\end{figure}

\subsubsection{Presión atmosférica}

La presión atmosférica de algunos sitios, medida con varios nodos dentro de Bogotá, se aprecia en la figura~\ref{gl7}. La medición más acertada se logra con los nodos Motell-0405 y Motell-0406. El nodo IM880b-i-01 presenta un desajuste, pero su medida mantiene la precisión esperada. El nodo mDot-Box-01 presenta varios errores de calibración.

\begin{figure}[H]
  \begin{center}
    \caption{Medidas múltiples de presión en sitios diversos}\label{gl7}
    \includegraphics[width =\textwidth]{Resultados/gl7.png}

    \propia
  \end{center}
\end{figure}

El nodo IM880 tomó las variaciones cíclicas diarias de la presión atmosférica en el décimo piso del edificio administrativo durante tres días completos. La figura \ref{gl8} presenta dicho ciclo, el cual se repite cada veinticuatro horas y consiste en el paso por dos mínimos: uno hacia las cuatro de la madrugada y otro hacia las cuatro de la tarde, y el paso por dos máximos: uno hacia las diez de la mañana y otro hacia las diez de la noche.

\begin{figure}[H]
  \begin{center}
    \caption{Medición local de presión en ciclos de veinticuatro horas}\label{gl8}
    \includegraphics[width=.6\textwidth]{Resultados/gl8.png}

    \propia
  \end{center}
\end{figure}

La diferencia entre los valores pico y valle de dicha magnitud está alrededor de 400 Pa, valor que se encuentra dentro de lo esperado en este tipo de fenómenos para la ciudad de Bogotá (Ideam, 2007).

\subsubsection{Temperatura}

En un periodo de quince días no continuos, se tomó la temperatura ambiente dentro de un espacio cerrado con cuatro nodos. Los nodos IM880b-01 y mDot-Box-01 presentaron fallas de ajuste en la medida, como lo muestra la figura~\ref{gl9}.

\begin{figure}[t]
  \begin{center}
    \caption{Mediciones de temperatura}\label{gl9}
    \includegraphics[width=.9\textwidth]{Resultados/gl9.png}

    \propia
  \end{center}
\end{figure}

Los nodos Motell-0405 y Motell-0406 presentaron el mejor comportamiento, con algunas pérdidas aleatorias de información.

\begin{figure}[b]
  \begin{center}
    \caption{Medición de temperatura dentro del laboratorio en ciclos de veinticuatro horas}
    \label{gl10}
    \includegraphics[width=.5\textwidth]{Resultados/gl10.png}

    \propia
  \end{center}
\end{figure}

La figura \ref{gl10} corresponde a la medición de temperatura realizada con el nodo IM880b-01 una vez corregido el error de calibración que presentó. El ambiente supervisado corresponde al mismo de la figura~\ref{gl9}.


\chapter*{Conclusiones}
\addcontentsline{toc}{chapter}{Conclusiones}

Los nodos implementados en el proyecto presentaron dificultades en la adaptación, dado que cada sensor entrega los datos según las condiciones del fabricante, por lo que se realizó el procesamiento de datos correspondiente para adecuarlos a los protocolos LoRaWAN.

El geolocalizador es uno de los sensores de mayor atención, en referencia a los utilizados en el nodo, dado que su sensibilidad permite determinar el lugar de la toma de datos. Si no hay geolocalización, el dato de interés se pierde, puesto que el nodo solo lo envía una o dos veces, de acuerdo con la clase de comunicación en la cual esté configurado.

El desempeño de este tipo de redes en la medición de parámetros provenientes de procesos, cuyo volumen de datos en el tiempo es bajo, queda demostrado con los levantamientos realizados durante las pruebas.

El nivel de pérdida de información es nulo, en parte por la pausa natural de los datos y, obviamente, por las técnicas aplicadas por el sistema. No se evidencian limitaciones respecto al tipo de variable ni al tipo de sensor empleado, en el contexto de monitoreo de parámetros ambientales.

La localización en interiores no cubrió las expectativas, puesto que LoRaWAN en clase A solo logró comunicaciones entre dos y tres pisos del edificio Sabio Caldas.

\enlargethispage{\baselineskip}

Los rangos de cobertura del sistema IoT LoRaWAN no alcanzan los valores indicados en las especificaciones del estándar, lo que puede atribuirse a que es una tecnología reciente. Las pruebas se realizaron con el sistema Multitech y con los equipos desarrollados en el GITUD.

La interoperabilidad entre los fabricantes (Multitech e ISMT) se comprobó mediante barridos entre la pasarela Multitech y el nodo ISMT, Mote II, el cual mostró una cobertura similar a la obtenida con el sistema Multitech.

El sistema LoRaWAN implementado en el grupo de investigación cumplió con las expectativas en cuanto a las comunicaciones y el funcionamiento de los nodos, logrando incluso una cobertura ligeramente superior a la del sistema Multitech, como se mostró en los mapas de cobertura.

Para futuros proyectos, se recomienda seleccionar sensores compatibles con los protocolos LoRaWAN, con el fin de minimizar el procesamiento de datos en el nodo. Asimismo, se recomienda emplear un geolocalizador de mayor sensibilidad para pruebas en interiores.

Con la tecnología LoRaWAN se pueden implementar sistemas IoT abiertos, dado que es compatible con \emph{software} y \emph{hardware} libres. Con la llegada de LoRaWAN clase C, es posible desarrollar sistemas inteligentes más robustos.

La ubicación de las pasarelas requiere de permisos institucionales, ya que, al tener antenas, deben instalarse en lugares elevados, con acceso a energía y conectividad. Para el caso de este proyecto, se consideró que la institución cuenta con varias sedes en Bogotá, lo que permitió realizar mediciones en dos puntos de la ciudad.

LoRaWAN clase A permite la configuración flexible de la red y, al ser un sistema abierto, puede adaptarse a diferentes tipos de aplicaciones y requisitos de la red, tanto libres como propietarios.

Los nodos LoRaWAN clase A solo pueden transmitir un número limitado de paquetes al gateway durante un periodo determinado. Esto puede representar una limitación para aplicaciones que requieren una alta tasa de transmisión de datos.

Si se requieren comunicaciones bidireccionales entre nodos y pasarelas, LoRaWAN ofrece las clases B y C, que permiten mayores tasas de datos que la clase A, aunque con una menor duración de la batería.


\cleardoublepage

\nocite{*}
\printbibliography

\chapter*{Siglas}
\thispagestyle{empty}

\begin{description}
  \item[μC]
    microcontrolador (\emph{microcontroller}).
  \item[ADC]
    conversor análogo a digital (\emph{analog to digital converter}).
  \item[ADR]
    tasa de datos adaptativa (\emph{adaptive data rate}).
  \item[AES]
    estándar de cifrado avanzado (\emph{advanced encryption standard}).
  \item[AQI]
    índice de calidad del aire (\emph{air quality index}).
  \item[BT]
    producto tiempo × ancho de banda (\emph{bandwidth-time product}).
  \item[CMAC]
    código de autenticación de mensajes basado en cifrado (\emph{cipher-based message authentication code}).
  \item[COV]
    compuestos orgánicos volátiles.
  \item[CRC]
    código de redundancia cíclica (\emph{cyclic redundancy code}).
  \item[CRT]
    cifrado en modo contador (\emph{counter mode encryption}).
  \item[CSS]
    espectro ensanchado \emph{chirp} (\emph{chirp spread spectrum}).
  \item[DevAddr]
    dirección del nodo (\emph{device address}).
  \item[EPA]
    Agencia de Protección Ambiental (Environmental Protection Agency).
  \item[ETSI]
    Instituto Europeo de Normas de Telecomunicaciones\foreignlanguage{english}{(European Telecommunications Standards Institute)}.
  \item[FCnt]
    contador de trama (\emph{frame counter}).
  \item[FCtrl]
    control de trama (\emph{frame control}).
  \item[FHDR]
    encabezado de trama (\emph{frame header}).
  \item[FHSS]
    espectro ensanchado por saltos de frecuencia (\emph{frequency-hopping spread spectrum}).
  \item[FOpts]
    opciones de trama (\emph{frame options}).
  \item[FPort]
    campo para el puerto (\emph{port field}).
  \item[FRMPayload]
    carga útil de la trama (\emph{frame payload}).
  \item[FSK]
    modulación por desplazamiento de frecuencia (\emph{frequency shift keying}).
  \item[GFSK]
    modulación por desplazamiento de frecuencia gaussiana (\emph{Gaussian frequency shift keying}).
  \item[GPIO]
    entradas y salidas de propósito general (\emph{general purpose input/output}).
  \item[HAL]
    capa de abstracción de \emph{hardware} (\emph{hardware abstraction layer}).
  \item[I\textsuperscript{2}C]
    Circuito Interintegrado (Inter-Integrated Circuit).
  \item[IBSG]
    Grupo de Soluciones de Negocios en Internet (Internet Business Solutions Group).
  \item[IoT]
    internet de las cosas (\emph{internet of things}).
  \item[ISM]
    industrial-científica-médica (\emph{industrial-scientific-medical}).
  \item[LAN]
    red de área local (\emph{local area network}).
  \item[LoRa]
    largo alcance (\emph{long range}).
  \item[LoRaWAN]
    red de área amplia LoRa (\emph{long range wide area network}).
  \item[LP-WAN]
    red de área amplia de baja potencia (\emph{low power wide area network}).
  \item[M2M]
    máquina a máquina (\emph{machine to machine}).
  \item[MAC]
    control de acceso al medio (\emph{media access control}).
  \item[MACPayload]
    carga útil de la capa MAC (\emph{MAC payload}).
  \item[MCU]
    unidad microcontroladora (\emph{microcontroller unit}).
  \item[MEMS]
    sistemas microelectromecánicos (\emph{microelectromechanical systems}).
  \item[MHDR]
    encabezado de la capa MAC (\emph{MAC header}).
  \item[MIC]
    código de integridad del mensaje (\emph{message integrity code}).
  \item[NAAQS]
    Estándares Nacionales de Calidad del Aire Ambiental (National Ambient Air Quality Standards).
  \item[NDIR]
    infrarrojo no dispersivo (\emph{non-dispersive infrared}).
  \item[PHDR]
    encabezado de la capa física (\emph{physical header}).
  \item[PHDR\_CRC]
    \textls[-20]{código de redundancia cíclica del encabezado (\emph{packet header CRC}).}
  \item[PHYPayload]
    carga útil de la capa física (\emph{physical payload}).
  \item[RFU]
    reservado para uso futuro (\emph{reserved for future use}).
  \item[R\textsubscript{s}]
    tasa de símbolos nominal (\emph{symbol rate}).
  \item[RTOS]
    sistema operativo en tiempo real (\emph{real-time operating system}).
  \item[SF]
    factor de ensanchamiento (\emph{spreading factor}).
  \item[SNR]
    relación señal a ruido (\emph{signal-to-noise ratio}).
  \item[SoC]
    sistema en un chip (\emph{system on a chip}).
  \item[SPI]
    interfaz periférica serial (\emph{serial peripheral interface}).
  \item[WAN]
    red de área amplia (\emph{wide area network}).
  \item[WSN]
    red inalámbrica de sensores (\emph{wireless sensor network}).
\end{description}

\printacronyms[title={Siglas},toctitle={Siglas}]

\newpage\thispagestyle{empty}

%% Editar Bio de autores %%
\pagestyle{plain}

\chapter*{\lhseries Sobre los autores}
\addcontentsline{toc}{chapter}{Sobre los autores}

\autorinfo{José Noé~Poveda}{\lipsum[1]}{jose.poveda@udistrital.edu.co}{https://orcid.org/0000-0001-8378-4215}

\autorinfo{Luis~Martín~Santamaría}{\lipsum[1]}{luis.santamaria@udistrital.edu.co}{https://orcid.org/0000-0001-8378-4215}

\autorinfo{Elvis~Eduardo Gaona García}{\lipsum[1]}{elvis.gaona@udistrital.edu.co}{https://orcid.org/0000-0001-8378-4215}

\autorinfo{Alex Francisco~Bustos}{\lipsum[1]}{alex.bustos@udistrital.edu.co}{https://orcid.org/0000-0001-8378-4215}
%%%%%%%%%%%%%%%%%%%%%%%%%%%

\newpage

\escolofon{octubre}{2025}

\end{document}
